\documentclass[11pt]{article}
\usepackage[margin=1in]{geometry}
\usepackage[T1]{fontenc}
\usepackage[utf8]{inputenc}
\usepackage[french]{babel}
\usepackage{amsmath,amssymb,amsthm}
\usepackage[colorlinks=true,linkcolor=blue,urlcolor=blue]{hyperref}
\newtheorem{problem}{Problème}
\newenvironment{refs}{\par\small\noindent\emph{Références :}\begin{itemize}\setlength\itemsep{0pt}}{\end{itemize}\normalsize}
\title{Hors de la Caverne \\ \Large Cryptographie et théorie de l'information}
\author{}
\date{}
\begin{document}
\maketitle
\noindent\emph{Des problèmes, pas de réponses.} Chaque problème ci-dessous est posé
sans solution. Les références indiquent où trouver les connaissances nécessaires.
\medskip


\section*{Lycée}

\begin{problem}[{Cassez ce chiffre de César --- Lycée}]
\textbf{CRY-01.} \textbf{Problème.} Le message suivant a été chiffré par un chiffre
de César : chaque lettre du texte clair, rédigé en anglais, a été décalée vers l'avant dans
l'alphabet d'une même quantité secrète (Z revenant à A), les espaces restant intacts.
 \texttt{DRO ZBSCYXOBC WSCDKUO DRO CRKNYGC YX DRO GKVV PYB DRO DRSXQC
DROWCOVFOC}
\begin{enumerate}
\item[(a)] Retrouvez le texte clair et le décalage \textit{sans essayer aveuglément les 25 décalages} :
comptez les fréquences des lettres dans le texte chiffré, comparez-les aux fréquences
connues de l'anglais (E, T, A, O en sont les lettres les plus courantes) et justifiez votre
identification.
\item[(b)] Expliquez en général pourquoi tout chiffre qui remplace toujours une même lettre claire
par une même lettre chiffrée doit laisser fuir le profil de fréquences du texte clair, si
compliquée que soit la règle de remplacement.
\end{enumerate}

\end{problem}

\noindent Suétone rapporte que Jules César utilisait la version à décalage de trois
de ce chiffre pour sa correspondance militaire. L'attaque systématique --- l'analyse
fréquentielle --- a d'abord été décrite par le philosophe arabe al-Kindī dans le Bagdad du
neuvième siècle, dans un manuscrit sur le déchiffrement des messages cryptographiques ;
c'est sans doute la naissance de la cryptanalyse, et elle a tué la substitution simple pour
toujours. Toute définition moderne de sécurité est, au fond, la promesse que ce genre de
fuite statistique ne peut pas se produire.

\begin{refs}
\item Contexte : Chiffrement par décalage (chiffre de César) --- Wikipédia (\url{https://fr.wikipedia.org/wiki/Chiffrement_par_d%C3%A9calage})
\item Contexte : Analyse fréquentielle --- Wikipédia (\url{https://fr.wikipedia.org/wiki/Analyse_fr%C3%A9quentielle})
\item Lecture : Simon Singh, \textit{The Code Book}, ch. 1
\item Lecture : Katz \& Lindell, \textit{Introduction to Modern Cryptography}, ch. 1
\end{refs}

\bigskip

\begin{problem}[{Échauffements d'arithmétique modulaire --- Lycée}]
\textbf{CRY-02.} \textbf{Problème.} Écrivons $a \equiv b \pmod{n}$ pour signifier que
$n$ divise $a-b$.
\begin{enumerate}
\item[(a)] Démontrez que les congruences s'additionnent et se multiplient : si
$a \equiv b$ et $c \equiv d \pmod{n}$, alors $a+c \equiv b+d$ et
$ac \equiv bd \pmod{n}$.
\item[(b)] Démontrez qu'un carré parfait est toujours $\equiv 0$
ou $1 \pmod 4$, et déduisez-en qu'aucun entier de la forme $4k+3$ n'est somme de deux
carrés.
\item[(c)] Trouvez, avec démonstration, les deux derniers chiffres de $3^{2026}$. (Indice sur la
méthode, non sur la réponse : les puissances de 3 modulo 100 doivent finir par se répéter ---
pourquoi ?)
\end{enumerate}

\end{problem}

\noindent La notation des congruences est due à Gauss, qui ouvrit avec elle ses
\textit{Disquisitiones Arithmeticae} (1801) et transforma l'\og{} arithmétique de l'horloge \fg{} en
un calcul. Toute la cryptographie à clé publique --- Diffie--Hellman, RSA, courbes elliptiques ---
est du calcul dans exactement cette arithmétique : ces exercices de doigté sont donc le droit
d'entrée de tout ce qui suit. La théorie des nombres qui les sous-tend est développée plus
avant sur la page Théorie des nombres (\url{pure-number-theory.html}).

\begin{refs}
\item Contexte : Arithmétique modulaire --- Wikipédia (\url{https://fr.wikipedia.org/wiki/Arithm%C3%A9tique_modulaire})
\item Lecture : Hoffstein, Pipher \& Silverman, \textit{An Introduction to Mathematical Cryptography}, ch. 1
\end{refs}

\bigskip

\begin{problem}[{Quels chiffres affines sont vraiment des chiffres ? --- Lycée}]
\textbf{CRY-03.} \textbf{Problème.} Numérotons les lettres $A=0,\dots,Z=25$. Un
\textit{chiffre affine} de clé $(a,b)$ chiffre la lettre $x$ en
$ax+b \pmod{26}$.
\begin{enumerate}
\item[(a)] Montrez que cette règle peut être défaite (de sorte que le déchiffrement soit
possible) exactement lorsque $\gcd(a,26)=1$.
\item[(b)] Montrez concrètement ce qui se gâte pour
$a=13$ : trouvez deux lettres différentes qui se chiffrent en la même lettre chiffrée.
\item[(c)] Comptez le nombre de clés $(a,b)$ valides.
\item[(d)] Expliquez pourquoi, malgré ses centaines
de clés au lieu de 25, le chiffre affine tombe exactement sous la même attaque par analyse
fréquentielle que le chiffre de César.
\end{enumerate}

\end{problem}

\noindent Le chiffre affine est le plus petit exemple d'un thème qui traverse toute
la cryptographie : une application de chiffrement doit être une bijection, et savoir si une
formule donne une bijection est une question de théorie des nombres (ici, l'inversibilité
modulo 26). La même condition de PGCD, portée à des modules de 2048 bits, décide quels
exposants sont des clés RSA légales.

\begin{refs}
\item Contexte : Chiffre affine --- Wikipédia (\url{https://fr.wikipedia.org/wiki/Chiffre_affine})
\item Lecture : Hoffstein, Pipher \& Silverman, \textit{An Introduction to Mathematical Cryptography}, ch. 1
\end{refs}

\bigskip

\begin{problem}[{Un espace de clés immense n'est pas la sécurité --- Lycée}]
\textbf{CRY-04.} \textbf{Problème.} Un chiffre par substitution général remplace chaque
lettre par une autre selon une permutation arbitraire de l'alphabet de 26 lettres : il y a
donc $26!$ clés.
\begin{enumerate}
\item[(a)] Démontrez que $26! > 10^{26}$.
\item[(b)] Montrez qu'un ordinateur
essayant un milliard ($10^9$) de clés par seconde mettrait plus de $10^9$ années
à les essayer toutes.
\item[(c)] Et pourtant, les chiffres par substitution appliqués à un texte
anglais de quelques centaines de lettres se cassent couramment à la main. Réconciliez cela
avec (b) : dites précisément quelle propriété du chiffre l'attaque par analyse fréquentielle
exploite, et formulez une leçon de la forme \og{} un grand espace de clés est une condition
\_\_\_\_\_\_ de sécurité, mais non une condition \_\_\_\_\_\_ \fg{}.
\end{enumerate}

\end{problem}

\noindent C'est la plus ancienne leçon chèrement acquise du domaine : la recherche
exhaustive n'est pas la seule attaque. Edgar Allan Poe a mis en scène le cassage à la main
de chiffres par substitution dans \og{} Le Scarabée d'or \fg{} (1843), et al-Kindī possédait la
méthode un millénaire plus tôt. La cryptographie moderne répond par des définitions : un
schéma doit résister à \textit{toute} attaque efficace, et non seulement à l'attaque évidente ---
raison pour laquelle ce sont les démonstrations de sécurité, et non les décomptes de clés,
qui font la monnaie du domaine.

\begin{refs}
\item Contexte : Chiffrement par substitution --- Wikipédia (\url{https://fr.wikipedia.org/wiki/Chiffrement_par_substitution})
\item Contexte : Attaque par force brute --- Wikipédia (\url{https://fr.wikipedia.org/wiki/Attaque_par_force_brute})
\item Lecture : Katz \& Lindell, \textit{Introduction to Modern Cryptography}, ch. 1
\end{refs}

\bigskip

\begin{problem}[{Le chiffre de Vigenère et l'examen de Kasiski --- Lycée}]
\textbf{CRY-24.} \textbf{Problème.} Le chiffre de Vigenère fixe un mot-clé, disons \texttt{LEMON}, et
chiffre un message en décalant ses 1re, 2e, 3e, \dots{} lettres
par les 1re, 2e, 3e, \dots{} lettres du mot-clé, en répétant
celui-ci autant de fois qu'il le faut (avec $A = 0, \dots, Z = 25$ comme quantités de
décalage).
\begin{enumerate}
\item[(a)] Expliquez précisément pourquoi l'analyse fréquentielle à un seul comptage de CRY-01
échoue contre ce chiffre : une même lettre claire n'est plus toujours chiffrée par la même
lettre chiffrée.
\item[(b)] Montrez cependant que si le mot-clé a pour longueur $k$, alors les lettres du texte
chiffré aux positions $1, k+1, 2k+1, \dots$ sont toutes chiffrées par un seul et même
décalage de César --- et de même pour chacune des $k-1$ autres classes de positions.
\item[(c)] Démontrez l'observation qui fonde l'examen de Kasiski : si un bloc de lettres claires se
répète, et si les deux occurrences commencent à des positions dont la différence est un
multiple de $k$, alors le texte chiffré se répète aussi. Expliquez comment le recueil des
distances entre blocs chiffrés répétés et la prise de leurs diviseurs communs exposent la
longueur $k$ du mot-clé.
\item[(d)] Assemblez l'attaque complète : devinez $k$, scindez le texte chiffré en $k$ chiffres
de César, et cassez chacun par analyse fréquentielle. Quelle propriété de la langue du texte
clair l'attaque entière consomme-t-elle ?
\end{enumerate}

\end{problem}

\noindent L'idée polyalphabétique remonte à Alberti (1467) et à Bellaso
(1553) ; elle fut attribuée à tort à Blaise de Vigenère et régna trois siècles durant
comme \textit{le chiffre indéchiffrable}. Charles Babbage le cassa en privé vers 1854, et
Friedrich Kasiski publia la méthode en 1863, faisant de la cryptanalyse une science
systématique. La leçon a survécu à son époque : un chiffre à clé courte répétée n'est qu'un
entrelacement de nombreux chiffres faibles, et le masque jetable de CRY-05 est exactement la
limite dans laquelle le \og{} mot-clé \fg{} ne se répète jamais.

\begin{refs}
\item Contexte : Chiffre de Vigenère --- Wikipédia (\url{https://fr.wikipedia.org/wiki/Chiffre_de_Vigen%C3%A8re})
\item Contexte : Cryptanalyse du chiffre de Vigenère (examen de Kasiski) --- Wikipédia (\url{https://fr.wikipedia.org/wiki/Cryptanalyse_du_chiffre_de_Vigen%C3%A8re})
\item Lecture : Simon Singh, \textit{The Code Book}, ch. 2
\end{refs}

\bigskip

\begin{problem}[{Le double chiffrement n'apporte rien --- Lycée, short}]
\textbf{CRY-25.} \textbf{Problème.} Démontrez que chiffrer par un décalage de César de
$a$ puis chiffrer le résultat par un décalage de $b$ revient exactement au décalage de
César unique de $a + b \bmod 26$ --- de sorte que le double chiffrement de César (ou le
centuple) n'est pas plus fort que le simple. Concluez que les 26 décalages forment un groupe
pour la composition, et expliquez en une phrase pourquoi cette propriété de clôture est
précisément ce qui rend ici le chiffrement répété sans valeur.
\end{problem}

\noindent Une famille de chiffres close par composition s'appelle un groupe, et pour
de telles familles itérer le chiffrement n'achète rien. Savoir si DES avait ce défaut fut une
question ouverte sérieuse jusqu'à ce que Campbell et Wiener démontrent en 1992 que DES n'est
\textit{pas} un groupe --- raison pour laquelle Triple-DES l'a véritablement renforcé.

\begin{refs}
\item Contexte : Chiffrement par décalage (chiffre de César) --- Wikipédia (\url{https://fr.wikipedia.org/wiki/Chiffrement_par_d%C3%A9calage})
\item Contexte : Groupe (mathématiques) --- Wikipédia (\url{https://fr.wikipedia.org/wiki/Groupe_(math%C3%A9matiques)})
\end{refs}

\bigskip

\begin{problem}[{Quelle est la force d'une phrase secrète ? --- Lycée, short}]
\textbf{CRY-26.} \textbf{Problème.} On construit une phrase secrète en choisissant quatre
mots indépendamment et uniformément au hasard dans une liste publique de $7776 = 6^5$ mots.
Montrez que le nombre de phrases secrètes possibles vaut environ $3{,}7 \times 10^{15}$ ---
soit à peu près $4\log_2 7776 \approx 51{,}7$ bits d'entropie --- et comparez-le à un mot de
passe de huit lettres minuscules aléatoires ($26^8$ possibilités, environ
$37{,}6$ bits) : lequel est le plus difficile à deviner, et par quel facteur ?
\end{problem}

\noindent C'est le procédé Diceware (Arnold Reinhold, 1995) : la liste compte
$6^5$ mots afin que cinq jets de dés en sélectionnent un uniformément. L'idée, popularisée
par une bande dessinée d'xkcd, est celle de Shannon (CRY-11) déguisée --- ce qui compte est
l'entropie de la \textit{procédure de sélection}, non l'air compliqué du résultat, et le
caractère public de la procédure ne coûte rien.

\begin{refs}
\item Contexte : Diceware --- Wikipédia (\url{https://fr.wikipedia.org/wiki/Diceware})
\item Contexte : Robustesse d'un mot de passe --- Wikipédia (\url{https://fr.wikipedia.org/wiki/Robustesse_d%27un_mot_de_passe})
\end{refs}

\bigskip

\begin{problem}[{Clés de contrôle : attraper les erreurs que les gens font vraiment --- Lycée}]
\textbf{CRY-36.} \textbf{Problème.} Un ISBN-10 est une suite de chiffres $d_1 d_2 \dots d_{10}$, dont le
dernier peut aussi être le symbole X valant 10, vérifiant
\[ \sum_{i=1}^{10} i \, d_i \;\equiv\; 0 \pmod{11}. \]
Le dernier chiffre est choisi pour que cela soit vérifié.
\begin{enumerate}
\item[(a)] Démontrez que ce contrôle détecte toute erreur modifiant un seul chiffre, et toute erreur
échangeant deux chiffres distincts --- qu'ils soient adjacents ou non. Où exactement votre
démonstration utilise-t-elle que 11 est premier ?
\item[(b)] Le contrôle des cartes bancaires (Luhn) travaille modulo 10 : en lisant depuis la droite,
on double un chiffre sur deux, on retranche 9 de toute valeur doublée dépassant 9, et l'on
exige que le total soit $0 \bmod 10$. Montrez qu'il attrape toute erreur d'un chiffre et
toute transposition de chiffres \textit{adjacents} sauf une, et trouvez cette exception.
\item[(c)] Démontrez que l'exception est inévitable pour les schémas de cette forme : aucun contrôle
du type $\sum_i a_i d_i \equiv 0 \pmod{10}$, à poids entiers fixés $a_i$, ne peut
détecter à la fois toutes les erreurs d'un chiffre et toutes les transpositions adjacentes.
(Montrez que la détection des erreurs simples force tous les $a_i$ à être impairs, et voyez
ce que cela fait à $a_i - a_{i+1}$.)
\end{enumerate}

\end{problem}

\noindent En 1969 Jacobus Verhoeff fit une chose inhabituelle pour un mathématicien :
avant de concevoir un code, il collecta des données sur les fautes que commettent les êtres
humains en recopiant des nombres. Les chiffres isolément faux représentaient environ quatre
erreurs sur cinq, les transpositions adjacentes l'essentiel du reste --- ce sont donc les deux
classes qu'une clé de contrôle décimale doit attraper, et la question (c) montre que
l'arithmétique modulo 10 ne peut attraper les deux. La réponse de Verhoeff fut d'abandonner
tout à fait l'addition et de calculer dans le groupe diédral d'ordre 10, dont la
non-commutativité fournit exactement ce qui manque à l'addition modulaire. L'autre solution,
celle des libraires, fut de changer de module : l'ISBN-10 a pris 11, au prix d'un onzième
symbole X. Lorsque le commerce du livre a fusionné avec le système mondial de codes-barres en
2007, l'ISBN-13 a adopté les poids 1 et 3 modulo 10 et a discrètement abandonné la garantie
sur les transpositions.

\begin{refs}
\item Contexte : Clé de contrôle --- Wikipédia (\url{https://fr.wikipedia.org/wiki/Cl%C3%A9_de_contr%C3%B4le})
\item Contexte : Formule de Luhn --- Wikipédia (\url{https://fr.wikipedia.org/wiki/Formule_de_Luhn})
\item Contexte : International Standard Book Number --- Wikipédia (\url{https://fr.wikipedia.org/wiki/International_Standard_Book_Number})
\item Lecture : J. Verhoeff, \textit{Error Detecting Decimal Codes}, Mathematical Centre Tracts 29 (Amsterdam, 1969)
\end{refs}

\bigskip

\begin{problem}[{Arbres de Merkle : un reçu de la taille d'un logarithme --- Lycée}]
\textbf{CRY-41.} \textbf{Problème.} Un journal public contient $n = 2^k$ enregistrements $x_1, \dots, x_n$.
Hachez chaque enregistrement pour obtenir les feuilles d'un arbre binaire, puis hachez à
répétition les paires voisines --- en remplaçant deux valeurs $u, v$ par $h(u \,\|\, v)$,
où $\|$ désigne la concaténation --- jusqu'à ce qu'il ne reste qu'une valeur : la
\textit{racine}. Supposez seulement que la fonction de hachage $h$ est \textit{résistante aux
collisions}, au sens où personne ne peut exhiber deux entrées différentes de même
sortie.
\begin{enumerate}
\item[(a)] Démontrez qu'il y a exactement $k = \log_2 n$ niveaux au-dessus des feuilles, et
comptez le nombre total de calculs de hachage nécessaires pour construire l'arbre entier.
Remarquez que la racine publiée est une unique valeur de hachage, si grand que soit $n$.
\item[(b)] Une \textit{preuve d'inclusion} pour l'enregistrement $x_i$ est la liste des valeurs
sœurs rencontrées sur le chemin de la feuille $i$ jusqu'à la racine. Montrez que cette
liste compte exactement $k$ entrées, et décrivez précisément le calcul par lequel un
vérificateur qui ne connaît que la racine et $x_i$ la contrôle.
\item[(c)] Démontrez la propriété de sécurité : si un exploitant de journal malhonnête convainc le
vérificateur qu'un certain $y \ne x_i$ occupe la position $i$ d'un arbre de \textit{même}
racine, alors quelque part le long des deux chemins il existe deux entrées différentes de
$h$ ayant la même sortie, contredisant la résistance aux collisions. (Comparez les deux
chemins niveau par niveau et considérez le niveau le plus bas où les valeurs coïncident.)
\item[(d)] Des chiffres, puis une subtilité. Avec $n = 2^{30}$ enregistrements et des valeurs de
hachage de 32 octets, calculez la taille en octets d'une preuve d'inclusion et comparez-la au
coût de l'envoi du journal entier. Expliquez ensuite ce qui se gâte si les feuilles et les
nœuds internes sont hachés par exactement la même règle : montrez qu'un attaquant pourrait
présenter un nœud interne comme s'il s'agissait d'un enregistrement, et dites ce qu'il faut
intégrer à l'entrée de $h$ pour l'empêcher.
\end{enumerate}

\end{problem}

\noindent Ralph Merkle a décrit les arbres de hachage dans sa thèse de doctorat à
Stanford en 1979 et les a brevetés la même année, comme moyen de signer de nombreux messages
avec une seule valeur publique ; l'idée a dormi jusqu'à ce que le monde des systèmes découvre
qu'elle résout le problème de prouver l'appartenance à une énorme collection à quelqu'un qui
n'est pas disposé à la télécharger. Git nomme chaque commit par un arbre de hachages, Bitcoin
résume les transactions de chaque bloc par une racine de Merkle, et Certificate Transparency ---
créé après que la compromission en 2011 de l'autorité de certification DigiNotar fut passée
inaperçue des semaines durant --- publie chaque certificat TLS émis dans un journal de Merkle en
ajout seul, de sorte qu'un navigateur peut exiger une preuve au lieu de se fier à une
promesse. La confusion décrite en (d) n'a rien d'hypothétique : la norme (RFC 6962) préfixe
les entrées de feuille par un octet nul et les entrées internes par un octet égal à un,
exactement pour cette raison.

\begin{refs}
\item Contexte : Arbre de Merkle --- Wikipédia (\url{https://fr.wikipedia.org/wiki/Arbre_de_Merkle})
\item Contexte : Certificate Transparency --- Wikipédia (en anglais) (\url{https://en.wikipedia.org/wiki/Certificate_Transparency})
\item Contexte : Fonction de hachage cryptographique --- Wikipédia (\url{https://fr.wikipedia.org/wiki/Fonction_de_hachage_cryptographique})
\item Lecture : Katz \& Lindell, \textit{Introduction to Modern Cryptography}, ch. 5
\end{refs}

\bigskip

\begin{problem}[{Pourquoi le mode ECB dessine le manchot --- Lycée, short}]
\textbf{CRY-42.} \textbf{Problème.} Un chiffrement par bloc chiffre des blocs de 16 octets
sous une clé $K$ par une bijection $E_K$ de l'ensemble de tous les blocs de 16 octets ; le
mode \textit{dictionnaire de codes électronique} (ECB) chiffre un long message en le découpant
en blocs $P_1, P_2, \dots$ et en envoyant $E_K(P_1), E_K(P_2), \dots$. Démontrez que
$P_i = P_j$ si et seulement si les $i$-ième et $j$-ième blocs chiffrés sont égaux --- pour
tout chiffrement par bloc, toute clé, et si robuste que soit $E_K$ --- et décrivez exactement
ce qu'une oreille indiscrète apprend, dès lors, du chiffrement ECB d'une image comportant de
larges plages d'une seule couleur.
\end{problem}

\noindent La démonstration standard est une image de Tux, le manchot de Linux,
chiffrée en mode ECB : les couleurs changent et le manchot reste parfaitement visible. ECB
fut l'un des quatre modes publiés avec DES en 1980 et il a survécu dans les logiciels depuis
lors, ce qui explique que cette image a probablement empêché plus de brèches réelles
qu'aucun théorème. L'acte d'accusation formel arrive plus loin sur cette page : un schéma ECB
est déterministe, et CRY-38(a) montre qu'aucun schéma déterministe ne peut satisfaire la
définition moderne de la sécurité, si incassable que soit le chiffrement sous-jacent.

\begin{refs}
\item Contexte : Mode d'opération (cryptographie) --- Wikipédia (\url{https://fr.wikipedia.org/wiki/Mode_d%27op%C3%A9ration_(cryptographie)})
\item Contexte : Chiffrement par bloc --- Wikipédia (\url{https://fr.wikipedia.org/wiki/Chiffrement_par_bloc})
\end{refs}

\bigskip


\section*{Licence}

\begin{problem}[{Le masque jetable a le secret parfait --- Licence}]
\textbf{CRY-05.} \textbf{Problème.} Messages, clés et textes chiffrés sont des chaînes
de $n$ bits ; la clé $K$ est tirée uniformément au hasard et utilisée une seule fois ; le
chiffrement est le OU exclusif bit à bit, $C = M \oplus K$. Shannon a appelé un schéma de
chiffrement \textit{parfaitement secret} si, pour toute loi a priori sur les messages, tout
message $m$ et tout texte chiffré $c$,
\[ \Pr[M=m \mid C=c] \;=\; \Pr[M=m], \]
c'est-à-dire si voir le texte chiffré ne change en rien les croyances de l'oreille
indiscrète.
\begin{enumerate}
\item[(a)] Démontrez que le masque jetable est parfaitement secret.
\item[(b)] Démontrez la réciproque de Shannon :
tout schéma parfaitement secret exige au moins autant de clés que de messages, de sorte que
le secret parfait ne peut jamais coûter moins cher que le masque jetable.
\end{enumerate}

\end{problem}

\noindent Le masque a été breveté par Gilbert Vernam (1919) et sa discipline
d'usage unique préconisée par Joseph Mauborgne, mais le \textit{théorème} --- à la fois qu'il est
incassable et que l'incassabilité a exactement ce prix --- est de Claude Shannon, dans
\og{} Communication Theory of Secrecy Systems \fg{} (1949), l'article qui a fondé la cryptographie
mathématique. La réciproque de (b) explique toute la forme du sujet moderne : puisque nous
refusons d'acheminer des clés gigantesques, nous nous contentons de la sécurité
calculatoire, et chaque schéma ci-dessous vit avec une hypothèse de difficulté non
démontrée.

\begin{refs}
\item Contexte : Masque jetable --- Wikipédia (\url{https://fr.wikipedia.org/wiki/Masque_jetable})
\item Contexte : Sécurité inconditionnelle --- Wikipédia (\url{https://fr.wikipedia.org/wiki/S%C3%A9curit%C3%A9_inconditionnelle})
\item Article : C. E. Shannon, \og{} Communication Theory of Secrecy Systems \fg{} (1949), à propos de l'article --- Wikipédia (en anglais) (\url{https://en.wikipedia.org/wiki/Communication_Theory_of_Secrecy_Systems})
\item Lecture : Katz \& Lindell, \textit{Introduction to Modern Cryptography}, ch. 2
\end{refs}

\bigskip

\begin{problem}[{Le masque utilisé deux fois : réemploi et ruine --- Licence}]
\textbf{CRY-06.} \textbf{Problème.} Deux messages ont été chiffrés avec le
\textit{même} masque jetable : $c_1 = m_1 \oplus k$ et $c_2 = m_2 \oplus k$.
\begin{enumerate}
\item[(a)] Démontrez que $c_1 \oplus c_2 = m_1 \oplus m_2$ : la clé disparaît entièrement.
\item[(b)] Expliquez pourquoi $m_1 \oplus m_2$, pour deux messages en langue naturelle, est loin
d'être aléatoire, et décrivez une méthode pour en récupérer les deux messages.
\item[(c)] Mettez votre
méthode à l'œuvre. Les deux textes clairs ci-dessous sont des phrases anglaises en
majuscules et espaces (octets ASCII), chiffrées avec le même masque de 33 octets ; textes
chiffrés en hexadécimal : 
\texttt{c1 = 6dc1c864836e81da216d2b5af91cd46c12d1c6578a41f1c8b52428baebbe9f304b} 
\texttt{c2 = 74ccc810f36297a93776594a9110a7030dc2a9479c4efcdcb8243da185b9973751} 
Retrouvez les deux phrases.
\item[(d)] Réconciliez cette rupture avec CRY-05 : quelle hypothèse
du théorème du secret parfait a été violée ?
\end{enumerate}

\end{problem}

\noindent Ce n'est pas une fable de salle de classe. Pendant et après la Seconde
Guerre mondiale, des chiffreurs soviétiques sous pression de production réutilisèrent des
pages de masque ; le projet Venona de l'armée américaine exploita exactement l'identité de
(a) pour lire des milliers de messages du KGB sur quatre décennies. La même bévue s'est
répétée dans les logiciels depuis lors --- un chiffrement de flux à nonce répété est un masque
utilisé deux fois --- d'où la maxime gravée au fronton du domaine : \og{} ne jamais réutiliser de
matière à clé \fg{}.

\begin{refs}
\item Contexte : Masque jetable --- Wikipédia (\url{https://fr.wikipedia.org/wiki/Masque_jetable})
\item Contexte : Projet Venona --- Wikipédia (\url{https://fr.wikipedia.org/wiki/Projet_Venona})
\item Lecture : Katz \& Lindell, \textit{Introduction to Modern Cryptography}, ch. 3
\end{refs}

\bigskip

\begin{problem}[{Exponentiation modulaire rapide --- Licence}]
\textbf{CRY-07.} \textbf{Problème.} La cryptographie à clé publique calcule sans cesse
$a^e \bmod n$ pour des nombres de plusieurs centaines de chiffres.
\begin{enumerate}
\item[(a)] Décrivez la
méthode \og{} élever au carré et multiplier \fg{} (carrés répétés pilotés par l'écriture binaire de
$e$) et démontrez qu'elle calcule $a^e \bmod n$ en au plus
$2\lfloor \log_2 e \rfloor$ multiplications modulo $n$.
\item[(b)] Expliquez pourquoi réduire
modulo $n$ après chaque multiplication est essentiel --- estimez le nombre de chiffres
qu'aurait $a^e$ lui-même pour $a$ et $e$ de 300 chiffres.
\item[(c)] Calculez
$5^{2026} \bmod 2027$ à la main, et montrez que la méthode naïve (2025 multiplications
successives) dépasse tout ordinateur lorsque l'exposant a 300 chiffres, tandis que
\og{} carré et multiplie \fg{} n'a besoin que de quelques milliers de multiplications.
\end{enumerate}

\end{problem}

\noindent Les carrés répétés sont anciens --- ils apparaissent dans la prosodie
sanskrite de Pingala (avant 200 av. J.-C.) pour calculer des puissances --- mais ils sont
devenus une infrastructure porteuse en 1976 : sans eux, Diffie--Hellman et RSA seraient
inutilisables, car leur sécurité exige des exposants si grands qu'une exponentiation naïve
survivrait à l'univers. C'est l'exemple standard d'une tâche d'apparence exponentielle qui
cache un algorithme logarithmique.

\begin{refs}
\item Contexte : Exponentiation modulaire --- Wikipédia (\url{https://fr.wikipedia.org/wiki/Exponentiation_modulaire})
\item Contexte : Exponentiation rapide --- Wikipédia (\url{https://fr.wikipedia.org/wiki/Exponentiation_rapide})
\item Lecture : Hoffstein, Pipher \& Silverman, \textit{An Introduction to Mathematical Cryptography}, ch. 1
\end{refs}

\bigskip

\begin{problem}[{Diffie--Hellman : correction et hypothèse exacte --- Licence}]
\textbf{CRY-08.} \textbf{Problème.} Fixez un grand nombre premier $p$ et un générateur $g$
du groupe multiplicatif modulo $p$. Alice choisit un secret $a$ et envoie
$A = g^a \bmod p$ ; Bob choisit un secret $b$ et envoie $B = g^b \bmod p$ ; chacun
élève ce qu'il a reçu à la puissance de son propre secret.
\begin{enumerate}
\item[(a)] Démontrez que les deux parties aboutissent
à la même valeur $g^{ab} \bmod p$.
\item[(b)] Soyez maintenant précis sur la sécurité. Définissez le
problème du logarithme discret (DLP : à partir de $g^a$ trouver $a$), le problème
Diffie--Hellman calculatoire (CDH : à partir de $g^a, g^b$ trouver $g^{ab}$) et le
problème décisionnel (DDH : distinguer $g^{ab}$ d'un élément aléatoire). Démontrez la
chaîne de réductions : un algorithme pour DLP résout CDH, et un algorithme pour CDH résout
DDH. Dites ensuite exactement lequel des trois doit être difficile pour que la clé espionnée
soit (i) irrécupérable, (ii) indistinguable de l'aléa --- et expliquez pourquoi \og{} DLP est
difficile \fg{} à lui seul n'est \textit{pas} connu comme suffisant.
\item[(c)] Montrez qu'un homme du milieu capable de modifier
le trafic met le protocole complètement en échec : les hypothèses de difficulté ne peuvent
donc pas remplacer l'authentification.
\end{enumerate}

\end{problem}

\noindent \og{} New Directions in Cryptography \fg{} (1976) de Whitfield Diffie et Martin
Hellman, s'appuyant sur les casse-tête de Ralph Merkle, proposait la première méthode
permettant à deux inconnus de convenir d'un secret en public --- l'acte fondateur de la
cryptographie à clé publique. (Ellis, Cocks et Williamson au GCHQ avaient trouvé les idées
plus tôt, dans le secret.) La question (b) est la vraie leçon du problème : la cryptographie
moderne vit ou meurt selon qu'elle énonce \textit{exactement} quel problème calculatoire
l'adversaire est supposé incapable de résoudre.

\begin{refs}
\item Contexte : Échange de clés Diffie--Hellman --- Wikipédia (\url{https://fr.wikipedia.org/wiki/%C3%89change_de_cl%C3%A9s_Diffie-Hellman})
\item Contexte : Hypothèse décisionnelle de Diffie--Hellman --- Wikipédia (\url{https://fr.wikipedia.org/wiki/Hypoth%C3%A8se_d%C3%A9cisionnelle_de_Diffie-Hellman})
\item Article : W. Diffie et M. Hellman, \og{} New Directions in Cryptography \fg{}, IEEE Transactions on Information Theory (1976)
\item Cours : Stanford, Dan Boneh, Cryptography I (Coursera) (\url{https://www.coursera.org/learn/crypto})
\end{refs}

\bigskip

\begin{problem}[{RSA : démontrez que le déchiffrement marche --- et que la version scolaire est cassée --- Licence}]
\textbf{CRY-09.} \textbf{Problème.} Soit $n = pq$ avec $p, q$ premiers distincts,
et choisissez $e$ avec $\gcd(e,(p-1)(q-1)) = 1$ et $d$ avec
$ed \equiv 1 \pmod{(p-1)(q-1)}$. Le chiffrement est $m \mapsto m^e \bmod n$, le
déchiffrement $c \mapsto c^d \bmod n$.
\begin{enumerate}
\item[(a)] Démontrez $m^{ed} \equiv m \pmod n$ pour
\textit{tout} $m \in \mathbb{Z}/n\mathbb{Z}$ --- d'abord lorsque $\gcd(m,n)=1$ à l'aide du
théorème d'Euler, puis en général via le petit théorème de Fermat appliqué séparément modulo
$p$ et modulo $q$ et le théorème des restes chinois.
\item[(b)] Montrez maintenant que le \og{} RSA scolaire \fg{}
tel qu'énoncé n'est pas un schéma de chiffrement sûr : il est déterministe, de sorte qu'un
adversaire capable de deviner une courte liste de textes clairs possibles (un vote :
OUI/NON) identifie le message en chiffrant les candidats ; et il est malléable --- à partir
d'un chiffré $c$ de $m$, produisez un chiffré valide de $2m$ sans connaître $m$.
\item[(c)] Expliquez
en un paragraphe à quoi sert un remplissage aléatoire.
\end{enumerate}

\end{problem}

\noindent Rivest, Shamir et Adleman ont trouvé le schéma en 1977, un an après que
Diffie--Hellman eut appelé de ses vœux une telle fonction \og{} à trappe \fg{} (Clifford Cocks
l'avait secrètement au GCHQ en 1973). La question (a) est un joyau de théorie des nombres
appliquée --- les pièces qui travaillent sont Fermat, Euler et le théorème des restes chinois,
développés comme mathématiques pures dans le groupe de problèmes sur la primalité de la page
Théorie des nombres (\url{pure-number-theory.html}), trois siècles avant que quiconque
puisse les dépenser. La question (b) enseigne la seconde leçon fondatrice : un schéma
correct n'est pas encore un schéma sûr ; la définition de la sécurité fait un vrai
travail.

\begin{refs}
\item Contexte : Chiffrement RSA --- Wikipédia (\url{https://fr.wikipedia.org/wiki/Chiffrement_RSA})
\item Contexte : Théorème d'Euler (arithmétique) --- Wikipédia (\url{https://fr.wikipedia.org/wiki/Th%C3%A9or%C3%A8me_d%27Euler_(arithm%C3%A9tique)})
\item Article : R. Rivest, A. Shamir, L. Adleman, \og{} A Method for Obtaining Digital Signatures and Public-Key Cryptosystems \fg{}, Communications of the ACM (1978)
\item Lecture : Hoffstein, Pipher \& Silverman, \textit{An Introduction to Mathematical Cryptography}, ch. 3
\end{refs}

\bigskip

\begin{problem}[{L'attaque des anniversaires : établissez la borne en $\surd$N --- Licence}]
\textbf{CRY-10.} \textbf{Problème.} Une fonction de hachage rend des valeurs dans un ensemble de taille
$N$ ; modélisez ses sorties sur des entrées distinctes comme des tirages uniformes
indépendants.
\begin{enumerate}
\item[(a)] Montrez
qu'après $k$ tirages la probabilité que toutes les valeurs soient distinctes vaut
$\prod_{i=1}^{k-1}\bigl(1 - \frac{i}{N}\bigr)$ et, à l'aide de $1-x \le e^{-x}$,
déduisez-en qu'une collision survient avec probabilité au moins
$1 - e^{-k(k-1)/2N}$.
\item[(b)] Concluez que $k \approx 1{,}18\sqrt{N}$ tirages suffisent
pour une collision de probabilité $1/2$ --- et démontrez une borne correspondante dans
l'autre sens, en montrant que $o(\sqrt N)$ tirages trouvent une collision avec probabilité
$o(1)$, de sorte que $\sqrt N$ est la vérité et non un artefact.
\item[(c)] Tirez-en la
conséquence pratique : une empreinte de $n$ bits n'offre qu'environ $n/2$ bits de
résistance aux collisions. Expliquez pourquoi SHA-256 a une sortie de 256 bits quand la
cible réelle est \og{} 128 bits de sécurité \fg{}.
\end{enumerate}

\end{problem}

\noindent C'est le paradoxe des anniversaires --- 23 personnes suffisent pour un
anniversaire commun --- transformé en arme. Les attaques par collision gouvernent les marges de
conception de toute fonction de hachage ; les ruptures de MD5 et SHA-1 (Wang et al.
2004-05 ; la collision SHAttered sur SHA-1, 2017) ont réussi en battant la borne générique
$\sqrt N$, ce qui est bien la raison pour laquelle cette borne doit devenir une seconde
nature. Le même $\sqrt N$ resurgit en CRY-13 et dans la recherche quantique de Grover --- une
étrange rime à travers le domaine.

\begin{refs}
\item Contexte : Attaque des anniversaires --- Wikipédia (\url{https://fr.wikipedia.org/wiki/Attaque_des_anniversaires})
\item Contexte : Paradoxe des anniversaires --- Wikipédia (\url{https://fr.wikipedia.org/wiki/Paradoxe_des_anniversaires})
\item Contexte : Fonction de hachage cryptographique --- Wikipédia (\url{https://fr.wikipedia.org/wiki/Fonction_de_hachage_cryptographique})
\item Lecture : Katz \& Lindell, \textit{Introduction to Modern Cryptography}, ch. 5
\end{refs}

\bigskip

\begin{problem}[{L'entropie à partir des premiers principes --- Licence}]
\textbf{CRY-11.} \textbf{Problème.} Pour une variable aléatoire $X$ prenant la valeur $i$
avec probabilité $p_i$, Shannon a défini l'entropie
\[ H(X) = -\sum_i p_i \log_2 p_i \qquad (\text{avec } 0\log 0 = 0). \]
Démontrez-en les propriétés de base à partir de rien.
\begin{enumerate}
\item[(a)] $H(X) \ge 0$, avec égalité si et seulement si
$X$ est déterministe.
\item[(b)] Pour un alphabet de taille $n$, $H(X) \le \log_2 n$,
avec égalité si et seulement si $X$ est uniforme (utilisez l'inégalité de Jensen, ou
démontrez à la main la convexité nécessaire).
\item[(c)] Si $X$ et $Y$ sont indépendantes alors
$H(X,Y) = H(X) + H(Y)$.
\item[(d)] Compréhension : Shannon a démontré que la continuité, la
croissance en $n$ pour les lois uniformes et une règle naturelle de regroupement (règle de
chaînage) forcent $H$ à valoir $-K\sum p_i \log p_i$ pour une constante $K$. Lisez
l'énoncé de ce théorème d'unicité et expliquez avec vos propres mots pourquoi l'axiome de
regroupement est celui qui porte tout le poids.
\end{enumerate}

\end{problem}

\noindent \og{} A Mathematical Theory of Communication \fg{} (1948) s'ouvre en demandant ce
que peut bien vouloir dire \og{} quantité d'information \fg{}, et répond par un théorème : les
axiomes n'admettent qu'une seule fonction. L'entropie s'est ensuite révélée être la réponse
exacte à deux questions concrètes d'ingénierie --- la compressibilité (CRY-19) et la capacité
d'un canal (CRY-20) --- ce qui est le schéma de toute la théorie : définir, caractériser, puis
faire correspondre la définition à une tâche opérationnelle. La même fonctionnelle
réapparaît en physique statistique ; voyez le problème sur la distribution de Boltzmann sur
la page Physique mathématique (\url{applied-physics.html}).

\begin{refs}
\item Contexte : Entropie de Shannon --- Wikipédia (\url{https://fr.wikipedia.org/wiki/Entropie_de_Shannon})
\item Article : C. E. Shannon, \og{} A Mathematical Theory of Communication \fg{} (1948), à propos de l'article --- Wikipédia (\url{https://fr.wikipedia.org/wiki/A_Mathematical_Theory_of_Communication})
\item Lecture : Cover \& Thomas, \textit{Elements of Information Theory}, ch. 2
\end{refs}

\bigskip

\begin{problem}[{L'inégalité de Kraft --- Licence}]
\textbf{CRY-12.} \textbf{Problème.} Un code binaire attribue à chacun de $n$ symboles
un mot de code, chaîne finie de bits ; c'est un \textit{code préfixe} si aucun mot de code
n'est préfixe d'un autre (le décodage n'a alors jamais besoin de regarder en avant).
\begin{enumerate}
\item[(a)] Démontrez l'inégalité de Kraft :
si un code préfixe a des longueurs de mots $\ell_1, \dots, \ell_n$, alors
\[ \sum_{i=1}^{n} 2^{-\ell_i} \le 1. \]
\item[(b)] Démontrez la réciproque : toutes longueurs vérifiant l'inégalité sont réalisées par un
code préfixe.
\item[(c)] Plus difficile (McMillan) : montrez que la même inégalité vaut pour tout
code uniquement déchiffrable, préfixe ou non --- les codes préfixes ne coûtent donc rien en
compression. (Considérez la puissance $k$-ième de $\sum 2^{-\ell_i}$ et dénombrez les
chaînes.)
\end{enumerate}

\end{problem}

\noindent L'inégalité de Kraft (1949, tirée de son mémoire de master au MIT ; la
réciproque de McMillan date de 1956) est le plancher combinatoire de toute la compression de
données : elle convertit \og{} déchiffrable \fg{} en une contrainte analytique sur les longueurs,
contre laquelle l'entropie peut être jouée comme borne inférieure. C'est la charnière entre
l'entropie abstraite de CRY-11 et le théorème de compression de CRY-19, et l'argument de
dénombrement sur les arbres de la question (a) est un petit classique de l'écriture des
démonstrations.

\begin{refs}
\item Contexte : Inégalité de Kraft--McMillan --- Wikipédia (\url{https://fr.wikipedia.org/wiki/In%C3%A9galit%C3%A9_de_Kraft})
\item Contexte : Codage de Huffman --- Wikipédia (\url{https://fr.wikipedia.org/wiki/Codage_de_Huffman})
\item Lecture : Cover \& Thomas, \textit{Elements of Information Theory}, ch. 5
\end{refs}

\bigskip

\begin{problem}[{Le chiffrement d'ElGamal : correction, malléabilité et hypothèse exacte --- Licence}]
\textbf{CRY-27.} \textbf{Problème.} Fixez un groupe cyclique $G$ d'ordre $q$ de générateur $g$.
La clé secrète d'Alice est $x$, sa clé publique $h = g^x$. Pour chiffrer
$m \in G$, Bob tire un $r$ uniforme frais et envoie le couple
\[ (c_1, c_2) = (g^r,\ m \cdot h^r). \]
\begin{enumerate}
\item[(a)] Démontrez la correction : Alice retrouve $m$ comme $c_2 \cdot (c_1^x)^{-1}$.
\item[(b)] Observez que le chiffrement est aléatoire --- un même $m$ donne des chiffrés différents
selon les jours --- de sorte que l'attaque \og{} deviner puis chiffrer \fg{} qui a cassé le RSA
scolaire en CRY-09(b) échoue. Montrez cependant qu'ElGamal est malléable : à partir d'un
chiffré de $m$, produisez un chiffré valide de $m \cdot m'$ sans connaître $m$.
\item[(c)] Démontrez que retrouver $m$ à partir d'un chiffré et de la clé publique est
\textit{équivalent} au problème Diffie--Hellman calculatoire de CRY-08 : un algorithme pour
l'une des deux tâches en fournit un pour l'autre.
\item[(d)] Dites pourquoi l'indistinguabilité des chiffrés équivaut à l'hypothèse décisionnelle de
Diffie--Hellman, et démontrez le sens le plus facile (celui qui casse DDH distingue les
chiffrés).
\end{enumerate}

\end{problem}

\noindent Taher ElGamal a publié le schéma en 1985 ; c'est l'échange de clés
Diffie--Hellman (CRY-08) figé en un schéma de chiffrement --- $h^r$ est exactement le secret
partagé, utilisé comme masque à usage unique. Ce fut le premier schéma à clé publique dont la
sécurité est \textit{équivalente} à une hypothèse standard nette, et non simplement impliquée
par elle, et ses descendants (sur les groupes de courbes elliptiques de CRY-14) tournent au
sein des protocoles modernes. La malléabilité de la question (b), défaut pour le chiffrement,
est la graine homomorphe des schémas de vote électronique.

\begin{refs}
\item Contexte : Cryptosystème de ElGamal --- Wikipédia (\url{https://fr.wikipedia.org/wiki/Cryptosyst%C3%A8me_de_ElGamal})
\item Article : T. ElGamal, \og{} A Public Key Cryptosystem and a Signature Scheme Based on Discrete Logarithms \fg{}, IEEE Transactions on Information Theory (1985)
\item Lecture : Hoffstein, Pipher \& Silverman, \textit{An Introduction to Mathematical Cryptography}, ch. 2
\end{refs}

\bigskip

\begin{problem}[{RSA avec le théorème des restes chinois --- et l'attaque par faute qu'il invite --- Licence}]
\textbf{CRY-28.} \textbf{Problème.} Soit $n = pq$ un module RSA d'exposant privé $d$ comme
en CRY-09. Les implémentations pratiques ne calculent jamais $c^d \bmod n$ directement.
\begin{enumerate}
\item[(a)] Démontrez le théorème des restes chinois pour deux modules premiers entre eux :
l'application $x \mapsto (x \bmod p,\ x \bmod q)$ est un isomorphisme d'anneaux
$\mathbb{Z}/n\mathbb{Z} \to \mathbb{Z}/p\mathbb{Z} \times \mathbb{Z}/q\mathbb{Z}$,
et donnez l'inverse explicite.
\item[(b)] Justifiez la procédure de déchiffrement par restes chinois : calculez
$m_p = c^{\,d \bmod (p-1)} \bmod p$ et $m_q = c^{\,d \bmod (q-1)} \bmod q$,
puis recombinez par le théorème. Démontrez que le résultat est égal à $c^d \bmod n$ (c'est
le petit théorème de Fermat qui autorise la réduction de l'exposant).
\item[(c)] Chiffrez l'économie. Avec \og{} carré et multiplie \fg{} (CRY-07), chaque exponentiation de
demi-taille emploie deux fois moins de multiplications, chacune sur des opérandes deux fois
plus courts. En supposant une multiplication scolaire (coût quadratique en la longueur des
opérandes), montrez que le facteur d'accélération total vaut environ 4.
\item[(d)] Et voici la piqûre (Boneh--DeMillo--Lipton, 1997). Un attaquant provoque un incident
matériel de sorte que $m_p$ soit calculé faussement tandis que $m_q$ est correct, et
obtient la sortie fautive $\tilde m$ ainsi qu'une signature correcte $m$ du même message.
Démontrez que $\gcd(\tilde m - m,\ n) = q$, de sorte qu'une seule faute factorise le
module. Quelle contre-mesure bon marché cela impose-t-il ?
\end{enumerate}

\end{problem}

\noindent L'astuce des restes chinois est aussi vieille que le déploiement de RSA et
figure dans toute implémentation sérieuse --- une accélération gratuite d'un facteur quatre
achetée au vieux théorème de Sunzi. L'attaque par faute de 1997 de Boneh, DeMillo et Lipton
a montré que la même structure est un couteau tourné vers l'intérieur : un seul bit basculé
pendant une moitié du calcul livre la clé. Elle a ouvert le champ moderne de l'analyse par
fautes et par canaux auxiliaires, où l'adversaire attaque la physique du calcul plutôt que
les mathématiques du schéma.

\begin{refs}
\item Contexte : Théorème des restes chinois --- Wikipédia (\url{https://fr.wikipedia.org/wiki/Th%C3%A9or%C3%A8me_des_restes_chinois})
\item Contexte : Chiffrement RSA --- Wikipédia (\url{https://fr.wikipedia.org/wiki/Chiffrement_RSA})
\item Article : D. Boneh, R. DeMillo, R. Lipton, \og{} On the Importance of Checking Cryptographic Protocols for Faults \fg{}, EUROCRYPT 1997
\end{refs}

\bigskip

\begin{problem}[{La borne de Hamming et le code parfait (7,4) --- Licence}]
\textbf{CRY-29.} \textbf{Problème.} Un code binaire de longueur $n$ est une partie
$C \subseteq \{0,1\}^n$ ; sa distance minimale est la plus petite distance de Hamming
(nombre de coordonnées qui diffèrent) entre mots de code distincts. Un code de distance
minimale $2e+1$ corrige $e$ erreurs de bits quelconques par décodage au mot de code le
plus proche.
\begin{enumerate}
\item[(a)] Démontrez la borne de Hamming (dite d'empilement de sphères) : si $C$ corrige $e$
erreurs, alors
\[ |C| \cdot \sum_{i=0}^{e} \binom{n}{i} \;\le\; 2^n . \]
Un code atteignant l'égalité est dit \textit{parfait} : les boules de rayon $e$ centrées sur
les mots de code pavent exactement $\{0,1\}^n$.
\item[(b)] Construisez le code de Hamming $(7,4)$ : les 16 mots binaires
$(d_1,d_2,d_3,d_4,p_1,p_2,p_3)$ dont les trois bits de parité vérifient
$p_1 = d_1{+}d_2{+}d_4$, $p_2 = d_1{+}d_3{+}d_4$, $p_3 = d_2{+}d_3{+}d_4$
(mod 2). Démontrez que sa distance minimale vaut 3 et vérifiez qu'il atteint la borne de (a)
avec $e = 1$ : il est parfait.
\item[(c)] Exhibez concrètement la correction d'une erreur : montrez que les trois contrôles de
parité, convenablement lus, nomment la position de tout bit basculé (le syndrome \textit{est}
l'adresse de l'erreur en binaire).
\item[(d)] Compréhension : énoncez le théorème de classification (Tietäväinen 1973, après van
Lint) : les seuls codes binaires parfaits non triviaux sont les codes de Hamming et le
code de Golay $(23,12)$. Pourquoi la rareté de la perfection ne diminue-t-elle pas
l'utilité des codes imparfaits (rappelez-vous la promesse de CRY-20) ?
\end{enumerate}

\end{problem}

\noindent Richard Hamming, qui partageait un ordinateur à relais aux Bell Labs en
1947, s'est lassé de voir ses travaux du week-end mourir sur une unique erreur détectée mais
non corrigible, et il a inventé les premiers codes \textit{correcteurs} d'erreurs ; l'article
de 1950 a créé la théorie des codes. Le code $(7,4)$ est le plus petit cas intéressant de
la géométrie qui régit tout le sujet --- les mots de code comme centres de sphères dans
l'espace de Hamming --- et la borne d'empilement de sphères de (a) est l'ombre combinatoire
exacte de la limite de capacité de canal de Shannon (CRY-20).

\begin{refs}
\item Contexte : Borne de Hamming --- Wikipédia (\url{https://fr.wikipedia.org/wiki/Borne_de_Hamming})
\item Contexte : Code de Hamming (7,4) --- Wikipédia (\url{https://fr.wikipedia.org/wiki/Code_de_Hamming_(7,4)})
\item Article : R. W. Hamming, \og{} Error Detecting and Error Correcting Codes \fg{}, Bell System Technical Journal (1950)
\item Lecture : MacWilliams \& Sloane, \textit{The Theory of Error-Correcting Codes}, ch. 1
\end{refs}

\bigskip

\begin{problem}[{Compter les générateurs --- Licence, short}]
\textbf{CRY-30.} \textbf{Problème.} En admettant que le groupe multiplicatif
$(\mathbb{Z}/p\mathbb{Z})^\times$ est cyclique pour tout nombre premier $p$, démontrez
qu'il possède exactement $\varphi(p-1)$ générateurs, où $\varphi$ est l'indicatrice
d'Euler.
\end{problem}

\noindent Gauss a démontré la cyclicité dans les \textit{Disquisitiones} (1801) ;
le décompte importe en pratique parce que Diffie--Hellman (CRY-08) a besoin d'un générateur,
et ce résultat dit qu'une fraction honorable des éléments fera l'affaire --- l'une des raisons
pour lesquelles la génération de paramètres est faisable. La théorie des nombres se poursuit
sur la page Théorie des nombres (\url{pure-number-theory.html}).

\begin{refs}
\item Contexte : Racine primitive modulo n --- Wikipédia (\url{https://fr.wikipedia.org/wiki/Racine_primitive_modulo_n})
\item Lecture : Hoffstein, Pipher \& Silverman, \textit{An Introduction to Mathematical Cryptography}, ch. 1
\end{refs}

\bigskip

\begin{problem}[{Les fonctions ne peuvent pas créer d'entropie --- Licence, short}]
\textbf{CRY-31.} \textbf{Problème.} Avec l'entropie de CRY-11, démontrez que pour toute
fonction $f$ et toute variable aléatoire $X$,
$ H(f(X)) \le H(X) $,
avec égalité si et seulement si $f$ est injective sur le support de $X$.
\end{problem}

\noindent Ce fait d'apparence anodine est la raison pour laquelle aucune quantité de
hachage, d'étirement ou de \og{} brassage \fg{} ne peut améliorer un mot de passe faible : un
traitement déterministe ne peut que détruire de l'information, jamais en frapper. C'est le
cas zéro de l'inégalité de traitement de données, cheval de trait de la théorie de
l'information.

\begin{refs}
\item Contexte : Entropie de Shannon --- Wikipédia (\url{https://fr.wikipedia.org/wiki/Entropie_de_Shannon})
\item Lecture : Cover \& Thomas, \textit{Elements of Information Theory}, ch. 2
\end{refs}

\bigskip

\begin{problem}[{L'horloge comme oracle : les attaques temporelles comme modèle mathématique --- Licence}]
\textbf{CRY-37.} \textbf{Problème.} Un chiffre est une fonction ; une \textit{implémentation} est une fonction
plus un temps d'exécution, et ce temps d'exécution est une sortie supplémentaire que le
concepteur n'a jamais voulu publier. Modélisez l'exponentiation \og{} carré et multiplie \fg{}
(CRY-07) calculant $g^{d} \bmod N$ : en traitant les bits de $d$ depuis le haut, chaque
bit coûte une élévation au carré de coût $S$, et chaque bit égal à 1 coûte une
multiplication supplémentaire de coût $M$. Le temps mesuré vaut
$T = c + Sn + M\,w(d) + \xi$, où $n$ est la longueur en bits, $w(d)$ le poids de
Hamming, et $\xi$ un bruit de mesure de moyenne 0 et d'écart type
$\sigma$.
\begin{enumerate}
\item[(a)] Démontrez que $T$ est une fonction affine de $w(d)$, et estimez combien
d'échantillons de temps indépendants un adversaire doit prendre pour trancher entre deux
exposants candidats dont les poids de Hamming diffèrent de 1, avec une probabilité d'erreur
au plus $\delta$. Exprimez la réponse en fonction de $\sigma/M$ et de $\delta$, et
expliquez pourquoi le coût de l'attaque ne croît qu'avec le \textit{carré} du niveau de
bruit.
\item[(b)] Formalisez la récupération bit par bit de Kocher comme une procédure séquentielle : les
$j-1$ premiers bits de $d$ étant déjà connus, l'adversaire simule les deux
continuations sur chacune de nombreuses bases connues $g_1,\dots,g_m$ et compare les temps
prédits aux temps mesurés. Énoncez l'hypothèse de modélisation précise sous laquelle la
bonne devinette donne des temps résiduels de variance strictement plus petite que la
mauvaise, et démontrez-le sous cette hypothèse. Expliquez pourquoi cela fait passer la
récupération de clé d'exponentielle à linéaire en $n$.
\item[(c)] Le même phénomène sans aucune statistique : une étiquette d'authentification de message
de $n$ octets est vérifiée par une comparaison qui rend la main dès que deux octets
diffèrent. Démontrez qu'un adversaire capable de chronométrer le vérificateur retrouve
l'étiquette en $O(256n)$ requêtes plutôt que les $256^{n}$ qu'exige une recherche
exhaustive. Énoncez ensuite, comme propriété précise d'un programme, ce que doit vouloir dire
\og{} à temps constant \fg{} pour que cette attaque soit impossible, et vérifiez que l'échelle de
Montgomery pour l'exponentiation possède cette propriété alors que \og{} carré et multiplie \fg{} ne
l'a pas.
\item[(d)] Les contre-mesures aussi sont des théorèmes. L'aveuglement RSA calcule $m = (r^{e}x)^{d}
r^{-1} \bmod N$ pour un $r$ aléatoire frais. Démontrez que la valeur ainsi élevée à la
puissance est alors uniformément distribuée et indépendante de $x$, et déterminez
exactement lesquelles des attaques ci-dessus cela met en échec --- et lesquelles cela laisse
entièrement intactes.
\end{enumerate}

\end{problem}

\noindent Paul Kocher n'était pas un cryptographe universitaire lorsqu'il a remarqué,
en 1996, que les mathématiques de RSA ne disent rien du temps que prend RSA. L'article a
brisé une hypothèse véritable du domaine --- qu'on pouvait analyser la sécurité d'un chiffre à
partir de son seul comportement entrée-sortie --- et il a ouvert la cryptanalyse par canaux
auxiliaires, étendue par la suite à la consommation électrique, au rayonnement
électromagnétique, à l'acoustique et aux caches des processeurs. Brumley et Boneh ont montré
en 2003 que l'attaque fonctionne à travers un réseau contre OpenSSL, tuant l'espoir que le
bruit sauverait les systèmes réels ; la question (a) explique pourquoi, puisque le bruit ne
coûte à l'attaquant que davantage d'échantillons. La conséquence durable est un changement de
ce que doit dire une spécification : les normes cryptographiques modernes contraignent le
flot de contrôle et les motifs d'accès mémoire, et pas seulement l'algèbre, et \og{} à temps
constant \fg{} est désormais une propriété que l'on démontre.

\begin{refs}
\item Contexte : Attaque temporelle --- Wikipédia (\url{https://fr.wikipedia.org/wiki/Attaque_temporelle})
\item Contexte : Attaque par canal auxiliaire --- Wikipédia (\url{https://fr.wikipedia.org/wiki/Attaque_par_canal_auxiliaire})
\item Article : P. C. Kocher, \og{} Timing Attacks on Implementations of Diffie-Hellman, RSA, DSS, and Other Systems \fg{}, CRYPTO 1996
\item Article : D. Brumley \& D. Boneh, \og{} Remote Timing Attacks are Practical \fg{}, USENIX Security 2003
\end{refs}

\bigskip

\begin{problem}[{Registres à décalage à rétroaction linéaire : pourquoi la linéarité est fatale --- Licence}]
\textbf{CRY-43.} \textbf{Problème.} Un LFSR binaire de longueur $L$ détient un état
$(s_0, \dots, s_{L-1}) \in \mathbb{F}_2^L$ et le prolonge indéfiniment par la récurrence
\[ s_{j+L} \;=\; c_1 s_{j+L-1} + c_2 s_{j+L-2} + \cdots + c_L s_j \pmod 2 , \]
en produisant $s_0, s_1, s_2, \dots$. Un chiffrement de flux peut utiliser cette suite comme
suite chiffrante, en chiffrant par OU exclusif exactement comme le masque jetable de
CRY-05.
\begin{enumerate}
\item[(a)] Démontrez que la sortie est ultimement périodique de période au plus $2^L - 1$, et que
tout état initial non nul donne la période exactement $2^L - 1$ précisément lorsque le
polynôme de connexion $C(x) = 1 + c_1 x + \cdots + c_L x^L$ est primitif sur
$\mathbb{F}_2$. (Identifiez l'espace des états à $\mathbb{F}_2[x]/(C)$ et lisez la
récurrence comme la multiplication par $x$ ; lorsque $C$ est irréductible, c'est le corps
$\mathbb{F}_{2^L}$.)
\item[(b)] Cassez-le maintenant. Supposez qu'un adversaire apprenne $2L$ bits consécutifs de la
suite chiffrante --- ce qui arrive dès qu'un fragment du texte clair est deviné. Montrez que
les prises inconnues $c_1, \dots, c_L$ vérifient alors un système de $L$ équations
linéaires sur $\mathbb{F}_2$, et démontrez que la solution est unique lorsque la matrice
$L \times L$ associée est inversible. Concluez que tout le passé et tout l'avenir de la
suite chiffrante s'ensuivent, à un coût polynomial en $L$ plutôt que les $2^L$
qu'exigerait de deviner l'état.
\item[(c)] Berlekamp et Massey suppriment les deux défauts de (b) : aucune hypothèse
d'inversibilité, et $L$ n'a pas à être connu d'avance. Définissez la \textit{complexité
linéaire} d'une suite binaire finie comme la longueur du plus court LFSR qui la produit,
énoncez la garantie de l'algorithme --- il retrouve la complexité linéaire et les prises à
partir de $2L$ termes en $O(L^2)$ opérations sur les bits --- et démontrez le lemme
d'unicité sur lequel il repose : si deux LFSR de longueur au plus $L$ coïncident sur $2L$
termes consécutifs, ils engendrent la même suite infinie.
\item[(d)] Les concepteurs détruisent donc la linéarité plutôt que les registres. Pour chacun des
procédés suivants, expliquez quelle non-linéarité il introduit et par quoi l'attaque de (b)
doit être remplacée : combiner plusieurs LFSR par une fonction booléenne non linéaire ;
cadencer irrégulièrement un registre sous le contrôle d'un autre, comme dans l'A5/1 du GSM ;
filtrer un registre unique à travers une fonction de sortie non linéaire. Dans chaque cas,
dites quelle doit être la complexité linéaire de la suite obtenue pour que l'attaque de (b)
soit hors de portée.
\end{enumerate}

\end{problem}

\noindent Les registres à décalage ont été le cheval de trait des chiffrements de
flux matériels pendant un demi-siècle : quelques portes logiques, un bit par cycle d'horloge,
et une période démontrable. Solomon Golomb en a rassemblé les mathématiques dans \textit{Shift
Register Sequences} (1967), pour une bonne part développées au Jet Propulsion Laboratory
pour des signaux de télémétrie plutôt que pour le secret. Puis l'algorithme de décodage des
codes BCH d'Elwyn Berlekamp (1968) a été reformulé par James Massey en 1969 comme synthèse de
registres à décalage, et la complexité linéaire est devenue d'un même coup l'étalon standard
et la notice nécrologique standard des générateurs de suites chiffrantes --- une suite ne vaut
que le plus court registre capable de la produire. L'histoire des déploiements le confirme :
l'A5/1 du GSM, l'E0 de Bluetooth et le Content Scrambling System des DVD combinaient tous des
LFSR, et tous sont tombés.

\begin{refs}
\item Contexte : Registre à décalage à rétroaction linéaire --- Wikipédia (\url{https://fr.wikipedia.org/wiki/Registre_%C3%A0_d%C3%A9calage_%C3%A0_r%C3%A9troaction_lin%C3%A9aire})
\item Contexte : Algorithme de Berlekamp--Massey --- Wikipédia (en anglais) (\url{https://en.wikipedia.org/wiki/Berlekamp%E2%80%93Massey_algorithm})
\item Contexte : Chiffrement de flux --- Wikipédia (\url{https://fr.wikipedia.org/wiki/Chiffrement_de_flux})
\item Article : J. L. Massey, \og{} Shift-register synthesis and BCH decoding \fg{}, IEEE Transactions on Information Theory 15 (1969)
\end{refs}

\bigskip

\begin{problem}[{Le mode CBC et l'oracle de remplissage --- Licence}]
\textbf{CRY-44.} \textbf{Problème.} Le chaînage de blocs chiffrés répare le défaut de CRY-42. Avec des blocs
de $b$ octets, un chiffrement par bloc $E_K$ et un vecteur d'initialisation aléatoire
frais $C_0 = IV$, le chiffrement s'écrit
\[ C_i \;=\; E_K\bigl(P_i \oplus C_{i-1}\bigr), \qquad i = 1, \dots, m . \]
Les messages sont d'abord complétés jusqu'à un multiple de $b$ octets en ajoutant $t$
octets tous égaux au nombre $t$, où $1 \le t \le b$ (la règle PKCS\#7) ; le déchiffreur
contrôle que le remplissage retrouvé a bien cette forme.
\begin{enumerate}
\item[(a)] Démontrez que $P_i = E_K^{-1}(C_i) \oplus C_{i-1}$ inverse le chiffrement, et qu'avec un
$IV$ frais des blocs clairs égaux ne produisent plus des blocs chiffrés égaux. Montrez aussi
qu'altérer un bloc chiffré abîme exactement deux blocs clairs, et dites lesquels.
\item[(b)] Démontrez que CBC est malléable d'une manière très précise : basculer le bit $j$ de
$C_{i-1}$ bascule le bit $j$ de $P_i$ tout en transformant $P_{i-1}$ en une valeur
incontrôlée. Concluez que le chiffrement CBC n'offre absolument aucune garantie
d'intégrité.
\item[(c)] Passons à l'attaque. Supposez qu'un serveur déchiffre les textes chiffrés qu'on lui
soumet et ne révèle --- par un message d'erreur, une différence de temps ou tout autre canal ---
rien d'autre que le fait que le remplissage était bien formé ou non. Étant donné un couple
$(C_{i-1}, C_i)$, montrez comment un adversaire autorisé à soumettre des valeurs modifiées
de $C_{i-1}$ retrouve le dernier octet de $P_i$ en au plus 256 requêtes, puis l'octet qui
le précède, et ainsi de suite. Démontrez que le texte clair entier tombe en environ
$256\,b\,m$ requêtes, et énoncez avec soin l'ambiguïté qui surgit sur le tout premier octet
retrouvé et la façon dont l'adversaire la lève.
\item[(d)] Tirez la morale dans le vocabulaire de CRY-38. L'oracle de remplissage est une attaque à
textes chiffrés choisis, et le chiffrement CBC peut être IND-CPA sûr tout en y succombant
entièrement : IND-CPA n'est donc pas le mot de la fin. Définissez le chiffrement authentifié
et démontrez que la composition \og{} chiffrer puis authentifier \fg{} --- transmettre $C$ avec une
étiquette infalsifiable portant sur $C$, et refuser de déchiffrer tout ce dont l'étiquette
échoue --- rend l'attaque de (c) impossible.
\end{enumerate}

\end{problem}

\noindent Serge Vaudenay a présenté cette attaque à EUROCRYPT 2002 et, l'année
suivante, Canvel, Hiltgen, Vaudenay et Vuagnoux l'ont retournée contre de vrais serveurs
SSL/TLS par le seul chronométrage, alors que ces serveurs avaient été soigneusement écrits
pour rendre des messages d'erreur identiques. La même fissure s'est rouverte pendant une
décennie --- contre ASP.NET en 2010, contre la construction MAC-puis-chiffrement de TLS sous le
nom de \og{} Lucky Thirteen \fg{} en 2013, contre SSL 3.0 sous le nom de POODLE en 2014 --- parce qu'il
est vain de colmater la fuite tant que l'architecture y invite. TLS 1.3 a fini par supprimer
tout à fait les modes CBC et exige un chiffrement authentifié. La question (c) est la
démonstration préférée du domaine que \og{} l'adversaire apprend un bit \fg{} et \og{} l'adversaire
apprend tout \fg{} peuvent être la même phrase.

\begin{refs}
\item Contexte : Attaque par oracle de remplissage --- Wikipédia (en anglais) (\url{https://en.wikipedia.org/wiki/Padding_oracle_attack})
\item Contexte : Mode d'opération (cryptographie) --- Wikipédia (\url{https://fr.wikipedia.org/wiki/Mode_d%27op%C3%A9ration_(cryptographie)})
\item Contexte : Chiffrement authentifié --- Wikipédia (\url{https://fr.wikipedia.org/wiki/Chiffrement_authentifi%C3%A9})
\item Article : S. Vaudenay, \og{} Security Flaws Induced by CBC Padding --- Applications to SSL, IPSEC, WTLS\dots{} \fg{}, EUROCRYPT 2002
\end{refs}

\bigskip

\begin{problem}[{La distance d'unicité --- Licence, short}]
\textbf{CRY-45.} \textbf{Problème.} Shannon a défini la \textit{distance d'unicité} d'un chiffre
comme $n_0 = H(K)/D$, où $H(K)$ est l'entropie de la clé et
$D = \log_2 26 - H_{\mathrm{Eng}}$ la redondance par lettre de la langue du texte clair ;
au-delà de $n_0$ lettres de texte chiffré, on s'attend à ce qu'une seule clé donne un
déchiffrement doué de sens. En prenant $H_{\mathrm{Eng}} \approx 1{,}5$ bit par lettre pour
l'anglais et une clé uniforme, calculez $n_0$ pour le chiffre par substitution simple de
CRY-04 ($H(K) = \log_2 26!$), justifiez la formule en estimant le nombre espéré de clés
parasites qui survivent à $n$ lettres de texte chiffré, et expliquez pourquoi le masque
jetable de CRY-05 a $n_0 = \infty$.
\end{problem}

\noindent Shannon a introduit la distance d'unicité dans \og{} Communication Theory of
Secrecy Systems \fg{} (1949) et a mesuré l'entropie de l'anglais deux ans plus tard par une
expérience où des sujets humains devinaient la lettre suivante d'un texte, arrivant à environ
0,6 à 1,3 bit par lettre. Le nombre qui tombe ici --- une vingtaine de lettres --- explique
pourquoi les énigmes de CRY-01 et CRY-04 ont des réponses uniques : deux courtes lignes de
texte chiffré surdéterminent déjà la clé, et tout ce qui vient ensuite n'est que
confirmation. Un chiffre atteint une distance d'unicité infinie exactement lorsque la
redondance est nulle, ce qui fait de la compression avant chiffrement plus qu'une
économie.

\begin{refs}
\item Contexte : Distance d'unicité --- Wikipédia (\url{https://fr.wikipedia.org/wiki/Distance_d%27unicit%C3%A9})
\item Contexte : Entropie de Shannon --- Wikipédia (\url{https://fr.wikipedia.org/wiki/Entropie_de_Shannon})
\item Article : C. E. Shannon, \og{} Prediction and Entropy of Printed English \fg{}, Bell System Technical Journal 30 (1951)
\end{refs}

\bigskip


\section*{Master}

\begin{problem}[{Baby-step giant-step --- Master}]
\textbf{CRY-13.} \textbf{Problème.} Soit $G$ un groupe cyclique d'ordre connu
$n$ et de générateur $g$, et supposons $h = g^x$ pour un $x$ inconnu. L'algorithme de
Shanks pose $m = \lceil \sqrt n\, \rceil$, tabule les \og{} petits pas \fg{} $g^j$
pour $0 \le j < m$, puis calcule les \og{} grands pas \fg{} $h \cdot (g^{-m})^i$ pour
$i = 0, 1, 2, \dots$ jusqu'à trouver une entrée de la table.
\begin{enumerate}
\item[(a)] Démontrez que l'algorithme trouve toujours
$x$, et qu'il utilise $O(\sqrt n)$ opérations de groupe et $O(\sqrt n)$ éléments
stockés (précisez votre hypothèse sur les recherches en table).
\item[(b)] Comparez à la force brute
et exhibez le compromis temps--mémoire : pour une taille de table $m$ quelconque, quel est
le temps d'exécution ?
\item[(c)] Compréhension : Shoup a démontré (1997) que tout algorithme
\textit{générique} --- n'utilisant que les opérations de groupe, non le codage des éléments ---
requiert $\Omega(\sqrt p)$ opérations pour les groupes d'ordre premier $p$. Expliquez ce
que cela dit et ne dit pas des logarithmes discrets dans $(\mathbb{Z}/p\mathbb{Z})^\times$,
où les méthodes de calcul d'indices battent nettement $\sqrt p$.
\end{enumerate}

\end{problem}

\noindent Daniel Shanks a introduit la méthode (1971) pour des calculs de groupes de
classes, avant que quiconque sache que les logarithmes discrets garderaient le trafic du
monde. C'est la raison pour laquelle \og{} 128 bits de sécurité \fg{} à partir d'un groupe exige un
ordre d'environ $2^{256}$ : les attaques génériques (BSGS et le rho de Pollard) prennent la
racine carrée de la taille du groupe --- le même $\sqrt N$ qu'en CRY-10, pour la même raison
de collision. L'écart de la question (c) entre groupes génériques et
$(\mathbb{Z}/p\mathbb{Z})^\times$ est précisément ce qui permet aux groupes de courbes
elliptiques (CRY-14), où l'on ne connaît rien de mieux que les attaques génériques, d'employer
des clés bien plus courtes.

\begin{refs}
\item Contexte : Baby-step giant-step --- Wikipédia (\url{https://fr.wikipedia.org/wiki/Baby-step_giant-step})
\item Contexte : Logarithme discret --- Wikipédia (\url{https://fr.wikipedia.org/wiki/Logarithme_discret})
\item Lecture : Hoffstein, Pipher \& Silverman, \textit{An Introduction to Mathematical Cryptography}, ch. 2
\item Cours : MIT OCW 6.875 Cryptography and Cryptanalysis (\url{https://ocw.mit.edu/courses/6-875-cryptography-and-cryptanalysis-spring-2005/})
\end{refs}

\bigskip

\begin{problem}[{La loi de groupe des courbes elliptiques --- l'associativité est le point difficile --- Master}]
\textbf{CRY-14.} \textbf{Problème.} Soit $E$ la courbe $y^2 = x^3 + ax + b$
sur un corps de caractéristique $\ne 2, 3$, avec $4a^3 + 27b^2 \ne 0$, munie
d'un point $O$ \og{} à l'infini \fg{}. Définissez $P+Q$ par la règle de la corde et de la
tangente : la droite passant par $P$ et $Q$ rencontre $E$ en un troisième point ;
réfléchissez-le par rapport à l'axe des $x$.
\begin{enumerate}
\item[(a)] Établissez les formules d'addition explicites (les formules
à pente $\lambda$), cas du doublement compris.
\item[(b)] Vérifiez les axiomes de groupe faciles : $O$ est
neutre, les inverses sont les réfléchis, et la loi est commutative --- et expliquez
\textit{pourquoi} chacun de ces points est immédiat au vu de la géométrie.
\item[(c)] Le point difficile : démontrez
l'associativité, $(P+Q)+R = P+(Q+R)$. Toute voie honnête est recevable --- un calcul direct
(et brutal) avec les formules, l'argument classique par les neuf points d'intersection de
deux cubiques (Cayley--Bacharach), ou l'identification de $E$ à son groupe des classes de
diviseurs. Compréhension : expliquez pourquoi l'associativité est ici le seul axiome de
contenu réel, et ce que cela dit de la profondeur géométrique de la loi de groupe.
\end{enumerate}

\end{problem}

\noindent La construction de la corde et de la tangente remonte aux travaux de
Fermat et de Newton sur les équations diophantiennes, et la structure de groupe a été
cristallisée par Poincaré (1901). Koblitz et Miller ont proposé indépendamment (1985)
d'utiliser ces groupes en cryptographie, précisément parce que --- contrairement à
$(\mathbb{Z}/p\mathbb{Z})^\times$ --- on ne connaît aucune attaque par calcul d'indices, de
sorte que la barrière générique en $\sqrt N$ de CRY-13 semble être la vérité, et que des
courbes de 256 bits égalent un RSA de 3072 bits. Quiconque s'est une fois battu avec la
question (c) ne prend plus jamais \og{} ceci est un groupe \fg{} sur parole.

\begin{refs}
\item Contexte : Courbe elliptique --- Wikipédia (\url{https://fr.wikipedia.org/wiki/Courbe_elliptique})
\item Contexte : Cryptographie sur les courbes elliptiques --- Wikipédia (\url{https://fr.wikipedia.org/wiki/Cryptographie_sur_les_courbes_elliptiques})
\item Lecture : Silverman \& Tate, \textit{Rational Points on Elliptic Curves}, ch. 1
\item Lecture : Hoffstein, Pipher \& Silverman, \textit{An Introduction to Mathematical Cryptography}, ch. 6
\end{refs}

\bigskip

\begin{problem}[{L'algorithme de Shor : que calcule-t-il exactement ? --- Master}]
\textbf{CRY-15.} \textbf{Problème.} Ceci est un problème de compréhension avec un vrai
théorème à l'intérieur.
\begin{enumerate}
\item[(a)] Démontrez la réduction classique au cœur de l'algorithme de Shor :
soit $N$ impair ayant au moins deux facteurs premiers distincts, tirez un $a$ au hasard
premier avec $N$, et soit $r$ l'ordre multiplicatif de $a \bmod N$. Montrez
que si $r$ est pair et $a^{r/2} \not\equiv -1 \pmod N$, alors
$\gcd(a^{r/2} - 1,\, N)$ est un facteur non trivial de $N$. (Montrez aussi, ou bien
cherchez et énoncez, qu'un $a$ aléatoire réussit avec probabilité au moins $1/2$.)
\item[(b)] Expliquez ce que fait la partie quantique : elle trouve la période $r$ de la fonction
$k \mapsto a^k \bmod N$ à l'aide de la transformée de Fourier quantique, en temps polynomial
en $\log N$ --- la seule étape sans contrepartie classique connue.
\item[(c)] Assemblez la
conséquence : pourquoi un grand ordinateur quantique tolérant aux fautes casse-t-il RSA
(CRY-09), Diffie--Hellman (CRY-08) et la cryptographie sur les courbes elliptiques (CRY-14) ---
tout en ne faisant qu'affaiblir AES et les fonctions de hachage (Grover : accélération
quadratique seulement, donc doublement des tailles de clés) ? Soyez précis sur les problèmes
calculatoires que résout l'algorithme de Shor : la factorisation et le logarithme discret,
tous deux instances du problème du sous-groupe caché pour les groupes abéliens finis.
\end{enumerate}

\end{problem}

\noindent Peter Shor a annoncé l'algorithme en 1994 ; la version publiée date de
1997. Il a converti l'informatique quantique d'une curiosité en un projet d'échelle
civilisationnelle, et il est la raison entière de la migration post-quantique de CRY-21.
Notez ce qu'il ne dit \textit{pas} : aucune démonstration n'existe que la factorisation soit
classiquement difficile --- Shor a élargi l'écart entre ce que nous savons faire
quantiquement et ce que nous \textit{croyons} classiquement, ce qui est un énoncé sur notre
ignorance autant que sur notre savoir (voir CRY-23).

\begin{refs}
\item Contexte : Algorithme de Shor --- Wikipédia (\url{https://fr.wikipedia.org/wiki/Algorithme_de_Shor})
\item Contexte : Problème du sous-groupe caché --- Wikipédia (en anglais) (\url{https://en.wikipedia.org/wiki/Hidden_subgroup_problem})
\item Article : P. Shor, \og{} Polynomial-Time Algorithms for Prime Factorization and Discrete Logarithms on a Quantum Computer \fg{} (1997), arXiv:quant-ph/9508027 (\url{https://arxiv.org/abs/quant-ph/9508027})
\item Lecture : Nielsen \& Chuang, \textit{Quantum Computation and Quantum Information}, ch. 5
\end{refs}

\bigskip

\begin{problem}[{Réseaux et apprentissage avec erreurs : la porte d'entrée --- Master}]
\textbf{CRY-16.} \textbf{Problème.} Fixez un module $q$ et un vecteur secret
$s \in (\mathbb{Z}/q\mathbb{Z})^n$. Un échantillon LWE est un couple
$(a,\ \langle a, s\rangle + e \bmod q)$ où $a$ est uniforme et $e$ une petite erreur
aléatoire.
\begin{enumerate}
\item[(a)] Montrez qu'avec $e = 0$ le secret se retrouve à partir de $n$
échantillons indépendants par élimination de Gauss --- le bruit est donc l'unique source de
difficulté.
\item[(b)] Expliquez l'image réticulaire : retrouver $s$ à partir de $m$ échantillons est un
problème de décodage à distance bornée sur un réseau aléatoire de $\mathbb{Z}^m$.
\item[(c)] Construisez le schéma de chiffrement de Regev à partir de LWE (clé publique : des
échantillons LWE ; on chiffre un bit en sommant un sous-ensemble aléatoire et en ajoutant
$\text{bit} \cdot \lfloor q/2 \rfloor$) et démontrez que le déchiffrement est correct
pourvu que l'erreur accumulée reste sous $q/4$.
\item[(d)] Compréhension : énoncez la réduction du pire cas au cas moyen de Regev --- la difficulté de
LWE sur des instances \textit{aléatoires} découle de la difficulté dans le pire cas de problèmes
de réseau approchés (GapSVP/SIVP) --- et expliquez pourquoi une telle réduction est une forme de
preuve fondamentalement plus forte que tout ce que l'on connaît pour la factorisation.
\end{enumerate}

\end{problem}

\noindent Oded Regev a introduit LWE en 2005 ; en quinze ans, c'était devenu le
fondement des schémas que le NIST a normalisés contre l'attaque quantique (CRY-21) et du
premier chiffrement complètement homomorphe (Gentry 2009). Le lien pire cas / cas moyen de la
question (d), qui descend de la découverte d'Ajtai en 1996, en est le cœur conceptuel : la
cryptographie a besoin d'instances \textit{aléatoires} difficiles, et les réseaux sont le seul
endroit où les mathématiques en fournissent actuellement à la demande.

\begin{refs}
\item Contexte : Apprentissage avec erreurs --- Wikipédia (\url{https://fr.wikipedia.org/wiki/Apprentissage_avec_erreurs})
\item Contexte : Cryptographie fondée sur les réseaux --- Wikipédia (en anglais) (\url{https://en.wikipedia.org/wiki/Lattice-based_cryptography})
\item Article : O. Regev, \og{} On Lattices, Learning with Errors, Random Linear Codes, and Cryptography \fg{}, Journal of the ACM (2009 ; version de conférence STOC 2005)
\item Lecture : Katz \& Lindell, \textit{Introduction to Modern Cryptography} (chapitre post-quantique, 3e éd.)
\end{refs}

\bigskip

\begin{problem}[{Divulgation nulle de connaissance : le protocole de l'isomorphisme de graphes --- Master}]
\textbf{CRY-17.} \textbf{Problème.} Peggy affirme que deux graphes publics $G_0, G_1$
sont isomorphes et connaît un isomorphisme $\varphi$. Protocole, répété $t$
fois : Peggy choisit une permutation $\pi$ uniformément aléatoire et envoie
$H = \pi(G_0)$ ; Victor tire à pile ou face un bit $b \in \{0,1\}$ ; Peggy répond par une
permutation $\sigma$ qui doit vérifier $H = \sigma(G_b)$ ; Victor contrôle. Démontrez
les trois propriétés qui définissent le protocole.
\begin{enumerate}
\item[(a)] \textit{Complétude} --- une Peggy honnête convainc toujours
Victor (exhibez $\sigma$ pour les deux valeurs de $b$).
\item[(b)] \textit{Solidité}
--- si les graphes ne sont \textit{pas} isomorphes, tout prouveur tricheur est démasqué à chaque
tour avec probabilité au moins $1/2$, et passe donc les $t$ tours avec probabilité au plus
$2^{-t}$.
\item[(c)] \textit{Divulgation nulle parfaite} --- construisez un
simulateur qui, ne disposant que de $G_0, G_1$ et d'aucun isomorphisme, produit des
transcriptions de loi exactement identique à celle des vraies (rembobinage autorisé), et
expliquez soigneusement pourquoi cela formalise \og{} Victor n'a rien appris au-delà de la
véracité de l'affirmation \fg{}.
\item[(d)] Compréhension : expliquez pourquoi l'interaction et l'aléa sont tous deux
indispensables --- une version déterministe ou non interactive de ce protocole
s'effondre.
\end{enumerate}

\end{problem}

\noindent Les preuves à divulgation nulle de connaissance ont été inventées par
Goldwasser, Micali et Rackoff (1985), et ce protocole, dû à Goldreich, Micali et Wigderson
(1986-91), en est l'exemple canonique --- ils ont de plus montré que tout énoncé de NP admet
une preuve à divulgation nulle, à condition que des fonctions à sens unique existent. La
définition par simulateur de la question (c) est l'un des grands gestes définitionnels de
l'informatique : \og{} n'a rien appris \fg{} devient \og{} aurait pu engendrer la conversation tout seul \fg{}.
Quarante ans après, les preuves ZK vérifient des calculs à grande échelle dans des systèmes
déployés ; la définition, elle, n'a pas changé.

\begin{refs}
\item Contexte : Preuve à divulgation nulle de connaissance --- Wikipédia (\url{https://fr.wikipedia.org/wiki/Preuve_%C3%A0_divulgation_nulle_de_connaissance})
\item Contexte : Problème de l'isomorphisme de graphes --- Wikipédia (\url{https://fr.wikipedia.org/wiki/Probl%C3%A8me_de_l%27isomorphisme_de_graphes})
\item Article : S. Goldwasser, S. Micali, C. Rackoff, \og{} The Knowledge Complexity of Interactive Proof Systems \fg{}, SIAM Journal on Computing (1989)
\item Article : O. Goldreich, S. Micali, A. Wigderson, \og{} Proofs that Yield Nothing But Their Validity \fg{}, Journal of the ACM (1991)
\end{refs}

\bigskip

\begin{problem}[{Le partage de secret de Shamir : la sécurité à seuil, exactement --- Master}]
\textbf{CRY-18.} \textbf{Problème.} Pour partager un secret $s \in \mathbb{F}_p$
($p$ premier, $p > n$) entre $n$ participants de sorte que $t$ quelconques
puissent le reconstituer mais que $t-1$ quelconques n'apprennent rien : tirez un polynôme
$f$ uniformément aléatoire de degré $\le t-1$ avec $f(0) = s$, et donnez au participant
$i$ la part $(i, f(i))$.
\begin{enumerate}
\item[(a)] Démontrez la reconstitution : $t$ parts quelconques déterminent $f$ de façon unique
(interpolation de Lagrange --- démontrez l'unicité, et pas seulement la formule), et donc
déterminent $s$.
\item[(b)] Démontrez la sécurité à seuil parfaite : pour $t-1$ parts quelconques et
\textit{tout} secret candidat $s' \in \mathbb{F}_p$, le nombre de polynômes compatibles avec
les parts et avec $f(0) = s'$ est exactement le même --- de sorte que la loi conditionnelle du
secret sachant $t-1$ parts est égale à sa loi a priori, au sens de la théorie de
l'information de CRY-05, sans la moindre hypothèse de difficulté.
\item[(c)] Montrez que le schéma est linéaire : les parts de $s$ et de $s'$ s'additionnent en des
parts de $s + s'$, et expliquez en une phrase pourquoi cette propriété homomorphe importe
(elle est la graine du calcul multipartite sécurisé).
\end{enumerate}

\end{problem}

\noindent Adi Shamir a publié \og{} How to Share a Secret \fg{} en 1979 --- deux pages, une
idée : un polynôme de degré $t-1$ représente exactement $t$ nombres d'information.
(George Blakley a trouvé un schéma géométrique indépendamment la même année.) C'est la
sécurité parfaite au sens de la théorie de l'information dans la nature --- descendante
directe de l'étalon de Shannon de 1949 --- et, via la question (c), elle est devenue la graine
du calcul multipartite (Ben-Or--Goldwasser--Wigderson, 1988) : calculer sur des secrets
qu'aucune partie ne détient.

\begin{refs}
\item Contexte : Partage de clé secrète de Shamir --- Wikipédia (\url{https://fr.wikipedia.org/wiki/Partage_de_cl%C3%A9_secr%C3%A8te_de_Shamir})
\item Contexte : Interpolation lagrangienne --- Wikipédia (\url{https://fr.wikipedia.org/wiki/Interpolation_lagrangienne})
\item Article : A. Shamir, \og{} How to Share a Secret \fg{}, Communications of the ACM (1979)
\item Lecture : Katz \& Lindell, \textit{Introduction to Modern Cryptography} (chapitre sur le partage de secret)
\end{refs}

\bigskip

\begin{problem}[{Le théorème du codage de source --- Master}]
\textbf{CRY-19.} \textbf{Problème.} Soit $X$ un symbole aléatoire d'entropie
$H(X)$, et considérons les codes préfixes binaires pour $X$ de longueur moyenne
$L = \sum_i p_i \ell_i$.
\begin{enumerate}
\item[(a)] Démontrez la réciproque : tout code préfixe vérifie
$L \ge H(X)$. (Combinez l'inégalité de Kraft, CRY-12, avec la positivité de l'entropie
relative, ou inégalité de Gibbs --- que vous démontrerez.)
\item[(b)] Démontrez l'atteignabilité :
les longueurs de Shannon $\ell_i = \lceil -\log_2 p_i \rceil$ vérifient Kraft et donnent
$L < H(X) + 1$ ; concluez qu'en codant des blocs de $k$ symboles on amène le débit à
$1/k$ de l'entropie.
\item[(c)] Compréhension : pour des sources i.i.d., lisez l'énoncé de la
propriété d'équipartition asymptotique (des ensembles typiques de taille environ
$2^{nH}$) et expliquez comment elle transforme (a)+(b) en le théorème complet du codage de
source de Shannon --- une compression fiable est possible à tout débit supérieur à $H$ et
impossible en dessous de $H$. Expliquez avec vos propres mots pourquoi l'entropie, définie
par les axiomes de CRY-11, se révèle être exactement la limite opérationnelle de la
compression.
\end{enumerate}

\end{problem}

\noindent C'est le premier des deux piliers de l'article de Shannon de 1948 : il
donne au \og{} contenu en information \fg{} un sens d'ingénieur --- le nombre de bits que l'on ne peut
pas éviter de transmettre. Huffman a trouvé le code optimal par symbole en 1951, étudiant de
troisième cycle (dans le cours de Fano, pour échapper à un examen final), et chaque fichier
zip et chaque codec de diffusion depuis lors vivent dans le couloir $[H, H+1)$ que ce
problème construit.

\begin{refs}
\item Contexte : Théorème du codage de source de Shannon --- Wikipédia (\url{https://fr.wikipedia.org/wiki/Th%C3%A9or%C3%A8me_du_codage_de_source})
\item Contexte : Propriété d'équipartition asymptotique --- Wikipédia (en anglais) (\url{https://en.wikipedia.org/wiki/Asymptotic_equipartition_property})
\item Lecture : Cover \& Thomas, \textit{Elements of Information Theory}, ch. 3 et 5
\item Cours : MIT OCW 6.441 Information Theory (\url{https://ocw.mit.edu/courses/6-441-information-theory-spring-2016/})
\end{refs}

\bigskip

\begin{problem}[{Capacité du canal binaire symétrique --- Master}]
\textbf{CRY-20.} \textbf{Problème.} Le canal binaire symétrique $\mathrm{BSC}(p)$
transmet un bit et le bascule avec probabilité $p < 1/2$, indépendamment d'un usage à
l'autre. Shannon a défini la capacité d'un canal comme $C = \max I(X;Y)$, le maximum sur
les lois d'entrée de l'information mutuelle entre entrée et sortie.
\begin{enumerate}
\item[(a)] En partant des définitions ($I(X;Y) = H(Y) - H(Y \mid X)$), démontrez que la
capacité de $\mathrm{BSC}(p)$ vaut
\[ C = 1 - H_2(p), \qquad H_2(p) = -p\log_2 p - (1-p)\log_2(1-p), \]
en identifiant la loi d'entrée qui réalise le maximum, et vérifiez les valeurs extrêmes
$C = 1$ pour $p = 0$ et $C = 0$ pour $p = 1/2$, en expliquant les deux.
\item[(b)] Compréhension : lisez
l'énoncé du théorème du codage de canal bruité et expliquez précisément ce qu'il promet ---
pour tout débit $R < C$ il existe des codes de grande longueur de bloc à probabilité
d'erreur tendant vers 0, et pour $R > C$ il n'en existe pas --- et pourquoi cela était
stupéfiant en 1948, alors que les ingénieurs croyaient que la fiabilité exigeait de ralentir
indéfiniment.
\item[(c)] Expliquez le rôle du codage aléatoire dans l'argument de Shannon, et pourquoi
le théorème est fameusement non constructif.
\end{enumerate}

\end{problem}

\noindent C'est le second pilier de l'article de 1948 et le théorème fondateur de
l'âge numérique : le bruit ne plafonne pas la fiabilité, seulement le débit. L'écart de la
question (c) --- Shannon démontre que de bons codes existent sans en exhiber un seul --- a
alimenté soixante-dix ans de théorie des codes, de Hamming et Reed--Solomon aux turbocodes,
aux codes LDPC et aux codes polaires d'Arıkan (2009), la première famille explicite dont on
ait démontré qu'elle atteint la capacité du canal binaire symétrique.

\begin{refs}
\item Contexte : Canal binaire symétrique --- Wikipédia (\url{https://fr.wikipedia.org/wiki/Canal_binaire_sym%C3%A9trique})
\item Contexte : Théorème du codage de canal bruité --- Wikipédia (\url{https://fr.wikipedia.org/wiki/Th%C3%A9or%C3%A8me_du_codage_de_canal})
\item Lecture : Cover \& Thomas, \textit{Elements of Information Theory}, ch. 7
\item Cours : MIT OCW 6.441 Information Theory (\url{https://ocw.mit.edu/courses/6-441-information-theory-spring-2016/})
\end{refs}

\bigskip

\begin{problem}[{Le rho de Pollard pour les logarithmes discrets --- Master}]
\textbf{CRY-32.} \textbf{Problème.} Soit $G$ un groupe cyclique d'ordre premier $n$ et de générateur
$g$, et $h = g^x$ avec $x$ inconnu. La méthode rho de Pollard itère une
application $f: G \to G$ \og{} d'allure aléatoire \fg{} : on partitionne $G$ en trois classes à
peu près égales et l'on pose $f(y) = gy$, $hy$ ou $y^2$ selon la classe de
$y$, en partant de $y_0 = 1$.
\begin{enumerate}
\item[(a)] Montrez que le long de l'orbite $y_{i+1} = f(y_i)$ on peut maintenir des exposants
$(a_i, b_i)$ avec $y_i = g^{a_i} h^{b_i}$, et que toute collision
$y_i = y_j$ avec $b_i \not\equiv b_j \pmod n$ donne
\[ x \equiv (a_j - a_i)(b_i - b_j)^{-1} \pmod n . \]
\item[(b)] Démontrez que l'orbite de toute suite ultimement périodique a la forme d'un \og{} rho \fg{}
(queue plus cycle), et que l'astuce de détection de cycle de Floyd --- avancer une tortue d'un
pas et un lièvre de deux pas jusqu'à ce que $y_i = y_{2i}$ --- trouve une collision en
stockant $O(1)$ éléments et en un nombre de pas linéaire en la somme des longueurs de la
queue et du cycle.
\item[(c)] Analyse heuristique : en modélisant $f$ comme une fonction uniformément aléatoire,
servez-vous de la borne des anniversaires de CRY-10 pour justifier que le temps de collision
espéré est $O(\sqrt n)$. Concluez au marché : rho égale le temps d'exécution
$O(\sqrt n)$ de baby-step giant-step (CRY-13) tout en remplaçant sa mémoire en
$O(\sqrt n)$ par $O(1)$ --- et expliquez pourquoi cela fait de rho, et non de BSGS,
l'attaque générique pratique, celle à laquelle les clés des défis publics sur courbes
elliptiques ont effectivement succombé.
\end{enumerate}

\end{problem}

\noindent John Pollard a publié la méthode en 1978, en la nommant d'après la lettre
grecque à laquelle ressemble le dessin de son orbite. C'est la raison pour laquelle \og{} la
sécurité générique d'un groupe de 256 bits vaut 128 bits \fg{} est un énoncé sur des machines
réelles et non seulement sur des asymptotiques : c'est le rho parallélisé
(van Oorschot--Wiener) qui a résolu les problèmes de défi ECC publics, et son modèle de coût
étalonne toutes les tailles de clés elliptiques déployées. L'analyse rigoureuse de la marche
--- démontrer le $O(\sqrt n)$ heuristique pour le $f$ réel, non aléatoire --- a dû attendre
les travaux de Miller, Venkatesan et d'autres, des décennies plus tard.

\begin{refs}
\item Contexte : Algorithme rho de Pollard (logarithme discret) --- Wikipédia (\url{https://fr.wikipedia.org/wiki/Algorithme_rho_de_Pollard_(logarithme_discret)})
\item Contexte : Détection de cycle --- Wikipédia (\url{https://fr.wikipedia.org/wiki/D%C3%A9tection_de_cycle})
\item Article : J. M. Pollard, \og{} Monte Carlo Methods for Index Computation (mod p) \fg{}, Mathematics of Computation 32 (1978)
\item Lecture : Hoffstein, Pipher \& Silverman, \textit{An Introduction to Mathematical Cryptography}
\end{refs}

\bigskip

\begin{problem}[{Codage arithmétique : comprimer jusqu'à l'entropie, deux bits de surcoût en tout --- Master}]
\textbf{CRY-33.} \textbf{Problème.} Le codage de Huffman (CRY-19) gaspille jusqu'à un bit \textit{par symbole},
ce qui est ruineux quand un symbole a une probabilité proche de 1. Le codage arithmétique ne
dépense son surcoût qu'une fois par \textit{message}. Codez une chaîne $x_1 x_2 \cdots x_n$
sur un alphabet fini de probabilités de symboles connues par subdivisions emboîtées : partez
de $[0,1)$ et, à chaque étape, découpez l'intervalle courant en sous-intervalles de
longueurs proportionnelles aux probabilités des symboles, en descendant dans celui qui
correspond au symbole suivant.
\begin{enumerate}
\item[(a)] Démontrez que l'intervalle final a pour largeur exactement
$P(x_1)P(x_2)\cdots P(x_n)$, la probabilité du message entier, et que des messages
distincts de même longueur reçoivent des intervalles disjoints.
\item[(b)] Montrez que tout intervalle de largeur $w$ contient un rationnel dyadique spécifiable
par $\lceil -\log_2 w \rceil + 1$ bits, et concluez que le message peut être transmis avec
au plus
\[ -\log_2 P(x_1 \cdots x_n) + 2 \]
bits --- à deux bits près du contenu en information du message entier, de sorte que le débit
par symbole tend vers l'entropie pour les longs messages, sans aucune astuce de mise en
blocs.
\item[(c)] Calculez les deux schémas sur une source concrète : des bits i.i.d. avec
$P(0) = 0{,}99$. Montrez que Huffman reste coincé à 1 bit par symbole tandis que la borne de
(b) donne environ $0{,}081$ bit par symbole --- et réconciliez cela avec le théorème du codage
de source.
\item[(d)] Compréhension : expliquez comment le décodeur retrouve le message à partir du
nombre transmis (il rejoue les mêmes subdivisions), et pourquoi toute la construction est le
codage de Shannon--Fano--Elias étiré des symboles isolés aux messages entiers.
\end{enumerate}

\end{problem}

\noindent L'idée a circulé sans être publiée depuis Elias au début des années 1960 ;
Rissanen et Pasco l'ont rendue implémentable en 1976, et l'article de 1987 de Witten, Neal et
Cleary, accompagné de code fonctionnel, l'a rendue universelle. Des fourrés de brevets l'ont
tenue à l'écart des premiers formats (bzip2 l'a fameusement évitée), mais les brevets ont
expiré et ses descendants siègent aujourd'hui dans JPEG, dans le CABAC de H.264/H.265 et dans
les codeurs ANS de zstd --- toute vidéo diffusée en flux est un nombre dans un intervalle qui
rétrécit.

\begin{refs}
\item Contexte : Codage arithmétique --- Wikipédia (\url{https://fr.wikipedia.org/wiki/Codage_arithm%C3%A9tique})
\item Article : I. H. Witten, R. M. Neal, J. G. Cleary, \og{} Arithmetic Coding for Data Compression \fg{}, Communications of the ACM 30 (1987)
\item Lecture : MacKay, \textit{Information Theory, Inference, and Learning Algorithms}, ch. 6
\end{refs}

\bigskip

\begin{problem}[{Conditionner diminue l'entropie --- Master, short}]
\textbf{CRY-34.} \textbf{Problème.} Définissez l'entropie conditionnelle
$H(Y \mid X) = \sum_x p(x)\, H(Y \mid X = x)$ et démontrez que
$H(Y \mid X) \le H(Y)$, avec égalité si et seulement si $X$ et $Y$ sont
indépendantes.
\end{problem}

\noindent \og{} L'information ne nuit jamais \fg{} --- en moyenne. La démonstration doit
respecter une subtilité véritable : pour une observation \textit{particulière} $x$, l'entropie
$H(Y \mid X = x)$ peut dépasser $H(Y)$ ; seule l'espérance est garantie de baisser. Cette
inégalité est la positivité de l'information mutuelle déguisée, et le moteur de tout théorème
réciproque de cette page.

\begin{refs}
\item Contexte : Entropie conditionnelle --- Wikipédia (\url{https://fr.wikipedia.org/wiki/Entropie_conditionnelle})
\item Lecture : Cover \& Thomas, \textit{Elements of Information Theory}, ch. 2
\end{refs}

\bigskip

\begin{problem}[{Aucune mise en gage n'est à la fois parfaitement cachante et parfaitement liante --- Master, short}]
\textbf{CRY-35.} \textbf{Problème.} Un schéma de mise en gage permet à un émetteur de publier
$c = \mathrm{Com}(m, r)$ (message $m$, aléa $r$) puis de l'ouvrir plus tard ; il est
\textit{parfaitement cachant} si la loi de $c$ est identique pour tout $m$, et
\textit{parfaitement liant} si aucun $c$ ne peut être ouvert sur deux messages différents.
Démontrez qu'aucun schéma ne satisfait les deux propriétés.
\end{problem}

\noindent L'impossibilité force un choix, et les deux choix sont faits
quotidiennement : les mises en gage de Pedersen (1991) sont parfaitement cachantes mais
seulement calculatoirement liantes ; celles fondées sur le hachage font l'inverse. Une mise
en gage est l'enveloppe scellée du cryptographe, et ce petit théorème dit que toute enveloppe
doit se fier aux limites calculatoires de quelqu'un --- une répétition en deux lignes des
arguments d'impossibilité à la Shannon de CRY-05.

\begin{refs}
\item Contexte : Mise en gage --- Wikipédia (\url{https://fr.wikipedia.org/wiki/Mise_en_gage})
\item Lecture : Katz \& Lindell, \textit{Introduction to Modern Cryptography}
\end{refs}

\bigskip

\begin{problem}[{IND-CPA, écrit en toutes lettres --- Master}]
\textbf{CRY-38.} \textbf{Problème.} Fixez un schéma de chiffrement symétrique
$\Pi = (\mathrm{Gen}, \mathrm{Enc}, \mathrm{Dec})$ de paramètre de sécurité
$\lambda$. Dans l'expérience IND-CPA, le challenger tire
$k \leftarrow \mathrm{Gen}(1^{\lambda})$ et un bit $b$ ; l'adversaire $A$ interroge
librement un oracle de chiffrement $\mathrm{Enc}_k(\cdot)$, soumet deux messages
$m_0, m_1$ de même longueur, reçoit $c \leftarrow \mathrm{Enc}_k(m_b)$, peut interroger de
nouveau, et produit $b'$. Son avantage vaut
\[ \mathrm{Adv}(A) \;=\; \Bigl| \Pr[b' = b] - \tfrac{1}{2} \Bigr| . \]
\begin{enumerate}
\item[(a)] Écrivez la définition complètement et formellement --- le paramètre de sécurité, l'ordre
des quantificateurs, ce que veulent dire \og{} temps polynomial probabiliste \fg{} et
\og{} négligeable \fg{} --- de sorte que rien ne soit laissé à la bonne volonté du lecteur. Démontrez
ensuite qu'\textit{aucun} schéma déterministe et sans état n'est IND-CPA sûr, en exhibant un
adversaire d'avantage exactement $1/2$. Concluez que le RSA scolaire (CRY-09) est peu sûr
au sens de cette définition, alors même que personne ne sait factoriser son module.
\item[(b)] Démontrez que tout schéma IND-CPA sûr a des chiffrés strictement plus longs que ses
textes clairs : montrez que si $|\mathrm{Enc}_k(m)| = |m|$ pour tout $m$, alors la
correction du déchiffrement force $\mathrm{Enc}_k$ à être déterministe, puis appliquez
(a).
\item[(c)] Démontrez l'hybride multi-messages : si $\Pi$ est IND-CPA sûr pour un unique défi,
alors pour tout $q$ polynomial en $\lambda$ aucun adversaire ne distingue le chiffrement
d'un $q$-uplet $(m_1,\dots,m_q)$ de celui de $(m'_1,\dots,m'_q)$ avec un avantage
supérieur à $q$ fois l'avantage à un seul message. Repérez précisément où le facteur
$q$ est perdu et dites s'il est nécessaire.
\item[(d)] Montrez que la définition n'est pas le mot de la fin. Donnez un adversaire qui casse le
chiffrement en mode CBC dans le jeu à \textit{chiffrés} choisis tout en laissant intacte sa
sécurité IND-CPA, et énoncez la définition originale de sécurité sémantique de Goldwasser et
Micali, en esquissant pourquoi elle équivaut à IND-CPA.
\end{enumerate}

\end{problem}

\noindent Avant 1982, \og{} chiffrement sûr \fg{} était une description de ce qu'une attaque
devait échouer à faire, et chaque proposition tombait sous une attaque que nul n'avait pensé
à interdire. \textit{Probabilistic Encryption} de Goldwasser et Micali a remplacé la liste de
souhaits par un jeu --- un adversaire, un défi, un avantage --- et en a tiré deux conséquences
qu'aucun concepteur n'avait acceptées : le chiffrement doit être aléatoire, et les chiffrés
doivent être plus longs que les textes clairs. Tout ce qui figure sur cette page au-dessus de
la section Master est un schéma ; ce problème-ci est la grammaire dans laquelle on dit ce
qu'un schéma réalise, et mettre les quantificateurs dans le bon ordre est tout le sujet. Ce
travail a valu à Goldwasser et Micali le prix Turing 2012, et la question (d) montre que le
projet définitionnel ne s'est pas arrêté là : la sécurité CPA ne dit rien d'un adversaire
capable de soumettre des chiffrés au déchiffrement, ce qui est la façon dont les protocoles
réels sont attaqués.

\begin{refs}
\item Contexte : Indistinguabilité des chiffrés (niveaux d'attaques) --- Wikipédia (\url{https://fr.wikipedia.org/wiki/Niveaux_d%27attaques})
\item Contexte : Sécurité sémantique --- Wikipédia (\url{https://fr.wikipedia.org/wiki/S%C3%A9curit%C3%A9_s%C3%A9mantique})
\item Article : S. Goldwasser \& S. Micali, \og{} Probabilistic Encryption \fg{}, Journal of Computer and System Sciences 28 (1984)
\item Lecture : Katz \& Lindell, \textit{Introduction to Modern Cryptography}, ch. 3
\end{refs}

\bigskip

\begin{problem}[{Réseaux de Feistel et le théorème de Luby--Rackoff --- Master}]
\textbf{CRY-46.} \textbf{Problème.} Fixez $n$ et des fonctions $f_1, \dots, f_r \colon \{0,1\}^n \to
\{0,1\}^n$. Le réseau de Feistel à $r$ tours $\Psi(f_1, \dots, f_r)$ agit sur
$\{0,1\}^{2n}$ en itérant
\[ (L_{i-1},\, R_{i-1}) \;\longmapsto\; (L_i,\, R_i) \;=\; \bigl(R_{i-1},\;\; L_{i-1}
\oplus f_i(R_{i-1})\bigr). \]
\begin{enumerate}
\item[(a)] Démontrez que $\Psi(f_1, \dots, f_r)$ est une permutation de $\{0,1\}^{2n}$ pour
\textit{tout} choix des $f_i$, fussent-elles constantes, et écrivez-en l'inverse
explicitement. Expliquez pourquoi c'est là tout l'intérêt du schéma : on peut bâtir un
chiffre sur une fonction de tour que personne n'a besoin de savoir inverser.
\item[(b)] Montrez que deux tours ne suffisent jamais. Avec $f_1, f_2$ tirées uniformément au
hasard, exhibez un adversaire qui distingue $\Psi(f_1, f_2)$ d'une permutation uniformément
aléatoire de $\{0,1\}^{2n}$ au moyen de deux entrées choisies et d'un travail constant, avec
un avantage proche de 1.
\item[(c)] Énoncez précisément le théorème de Luby--Rackoff : pour $f_1, f_2, f_3$ indépendantes et
uniformément aléatoires, aucun adversaire faisant $q$ requêtes ne distingue
$\Psi(f_1,f_2,f_3)$ d'une permutation uniformément aléatoire de $\{0,1\}^{2n}$ avec un
avantage dépassant $O(q^2/2^n)$. Démontrez-le, ou au minimum démontrez le lemme central :
sauf si deux des requêtes entrent en collision sur une certaine valeur interne --- événement de
probabilité $O(q^2/2^n)$ par la borne des anniversaires de CRY-10 --- les réponses que voit
l'adversaire sont uniformes et indépendantes.
\item[(d)] Trois tours ne suffisent pas pour une permutation pseudo-aléatoire \textit{forte}, où
l'adversaire peut aussi interroger l'inverse. Exhibez un tel distingueur pour
$\Psi(f_1,f_2,f_3)$, énoncez le théorème à quatre tours qui répare le défaut, puis dites
soigneusement ce que ce théorème affirme et n'affirme pas au sujet de DES, dont les seize
fonctions de tour ne sont ni uniformément aléatoires ni indépendantes.
\end{enumerate}

\end{problem}

\noindent L'équipe de Horst Feistel chez IBM a construit le chiffre Lucifer vers 1971
et la structure est passée dans DES en 1977. La question (a) explique son emprise sur le
domaine pendant cinquante ans : l'inversibilité est gratuite, un concepteur peut donc dépenser
tout son budget à rendre la fonction de tour imprévisible. En 1988, Michael Luby et Charles
Rackoff ont transformé le savoir empirique en théorème --- trois tours d'une fonction
pseudo-aléatoire donnent une permutation pseudo-aléatoire, quatre en donnent une qui survit
aux attaques à chiffrés choisis --- produisant ainsi l'une des premières constructions dont la
sécurité est \textit{réduite} à celle de ses composants plutôt qu'argumentée par l'échec des
attaques connues. L'écart de la question (d) est l'endroit où la théorie s'arrête et où
l'ingénierie reprend : aucun chiffre déployé n'a de fonctions de tour aléatoires
indépendantes.

\begin{refs}
\item Contexte : Réseau de Feistel --- Wikipédia (\url{https://fr.wikipedia.org/wiki/R%C3%A9seau_de_Feistel})
\item Contexte : Permutation pseudo-aléatoire --- Wikipédia (en anglais) (\url{https://en.wikipedia.org/wiki/Pseudorandom_permutation})
\item Article : M. Luby et C. Rackoff, \og{} How to Construct Pseudorandom Permutations from Pseudorandom Functions \fg{}, SIAM Journal on Computing 17 (1988)
\item Lecture : Katz \& Lindell, \textit{Introduction to Modern Cryptography}, ch. 7
\end{refs}

\bigskip

\begin{problem}[{Cryptanalyse linéaire et différentielle d'un chiffre jouet --- Master}]
\textbf{CRY-47.} \textbf{Problème.} Soit $S \colon \{0,1\}^4 \to \{0,1\}^4$ une bijection (une boîte S) et
construisez un petit réseau de substitution-permutation sur des blocs de 16 bits : chaque tour
ajoute par OU exclusif une clé de tour, applique $S$ à chacun des quatre quartets, puis
permute les seize positions de bits par un câblage fixe. Faites quatre tours.
\begin{enumerate}
\item[(a)] Démontrez le \textit{lemme d'empilement} de Matsui : si $X_1, \dots, X_k$ sont des
variables aléatoires binaires indépendantes avec $\Pr[X_i = 0] = \tfrac12 + \varepsilon_i$,
alors
\[ \Pr[X_1 \oplus \cdots \oplus X_k = 0] \;=\; \tfrac12 \;+\; 2^{\,k-1}
\prod_{i=1}^{k} \varepsilon_i . \]
Déduisez-en que le biais d'une chaîne d'approximations s'effondre géométriquement avec sa
longueur, et expliquez pourquoi cela fait de \og{} quel est le biais de la meilleure approximation
sur un tour ? \fg{} la question de conception d'une boîte S.
\item[(b)] Bâtissez l'attaque linéaire. Définissez la table d'approximation linéaire de $S$ : pour
chaque couple de masques $(a,b)$, le biais de l'événement $a \cdot x \oplus b \cdot S(x) = 0$
pour $x$ uniforme. Calculez-la pour une boîte S de 4 bits explicite de votre choix,
enchaînez les meilleures approximations sur trois tours, et décrivez comment la relation
biaisée obtenue entre bits du texte clair, bits du chiffré et bits de clé permet de retrouver
la clé du dernier tour en classant les sous-clés candidates. Démontrez que distinguer un biais
$\varepsilon$ de l'absence de biais demande de l'ordre de $\varepsilon^{-2}$ textes clairs
connus, et repérez où se cache la constante.
\item[(c)] Bâtissez l'attaque différentielle. Définissez la table de distribution des différences de
$S$, dont l'entrée en $(\alpha, \beta)$ compte les $x$ tels que $S(x) \oplus S(x \oplus \alpha) =
\beta$. Démontrez que chaque ligne somme à 16 et que chaque entrée est paire, et expliquez
pourquoi une permutation uniformément aléatoire aurait presque toutes ses entrées proches de 1.
Construisez ensuite une caractéristique différentielle sur trois tours de probabilité $p$, et
montrez qu'environ $p^{-1}$ couples de textes clairs choisis suffisent à la détecter et à
trier les bons couples des mauvais.
\item[(d)] Compréhension : une boîte S résiste aux deux attaques exactement lorsque les deux tables
sont plates. Énoncez les critères de conception qui en découlent --- non-linéarité et uniformité
différentielle --- et expliquez comment la boîte S d'AES, l'inversion dans
$\mathbb{F}_{256}$ suivie d'une application affine sur $\mathbb{F}_2$, a été choisie de
sorte que sa probabilité différentielle maximale soit démontrablement $2^{-6}$ et son biais
linéaire maximal comparablement petit. Dites ensuite précisément ce que cela ne démontre
\textit{pas} au sujet d'AES.
\end{enumerate}

\end{problem}

\noindent Eli Biham et Adi Shamir ont introduit la cryptanalyse différentielle vers
1990 et l'ont employée contre le DES complet avec $2^{47}$ textes clairs choisis --- sur quoi
Don Coppersmith a révélé en 1994 que l'équipe d'origine d'IBM avait trouvé la même technique
en 1974 sous le nom d'\og{} attaque T \fg{} et avait réglé les boîtes S de DES pour y résister, la
raison en étant tenue classifiée deux décennies durant à la demande de la NSA. Un chiffre
conçu en 1974 contre une attaque publiée en 1990 est le plus étrange compliment de l'histoire
du domaine. La cryptanalyse linéaire de Mitsuru Matsui a suivi en 1993 et, en 1994, il a
retrouvé une clé DES avec $2^{43}$ textes clairs connus et cinquante jours de temps de
station de travail --- la première rupture expérimentale du chiffre complet. Ensemble, les deux
attaques ont converti la conception des chiffrements par bloc d'un art en un problème
d'optimisation à deux objectifs publiés, et tout chiffre sérieux depuis lors est livré avec
des bornes contre les deux.

\begin{refs}
\item Contexte : Cryptanalyse linéaire --- Wikipédia (\url{https://fr.wikipedia.org/wiki/Cryptanalyse_lin%C3%A9aire})
\item Contexte : Cryptanalyse différentielle --- Wikipédia (\url{https://fr.wikipedia.org/wiki/Cryptanalyse_diff%C3%A9rentielle})
\item Article : E. Biham et A. Shamir, \og{} Differential Cryptanalysis of DES-like Cryptosystems \fg{}, Journal of Cryptology 4 (1991) ; M. Matsui, \og{} Linear Cryptanalysis Method for DES Cipher \fg{}, EUROCRYPT 1993
\item Lecture : H. M. Heys, \og{} A Tutorial on Linear and Differential Cryptanalysis \fg{}, Cryptologia 26 (2002)
\end{refs}

\bigskip

\begin{problem}[{Le protocole de Schnorr : extraire le témoin --- Master, short}]
\textbf{CRY-48.} \textbf{Problème.} Dans un groupe cyclique $G = \langle g \rangle$ d'ordre
premier $q$, un prouveur qui connaît $x$ tel que $h = g^x$ envoie $a = g^r$ pour un
$r$ uniforme frais, reçoit un défi $c \in \mathbb{Z}/q\mathbb{Z}$, et répond par
$s = r + cx \bmod q$ ; le vérificateur accepte lorsque $g^s = a\,h^c$. Démontrez la
\textit{solidité spéciale} --- à partir de deux transcriptions acceptantes quelconques
$(a,c,s)$ et $(a,c',s')$ de même premier message et avec $c \ne c'$, le secret $x$ se
calcule --- et la \textit{divulgation nulle vis-à-vis d'un vérificateur honnête}, en exhibant un
simulateur qui produit des transcriptions de loi exactement correcte tout en ne connaissant
que $g$ et $h$.
\end{problem}

\noindent Claus-Peter Schnorr a publié ce protocole d'identification et le schéma de
signature qui en dérive entre 1989 et 1991, en le concevant assez court pour tenir dans une
carte à puce. C'est le protocole canonique à trois passes, et les deux propriétés démontrées
ici en sont le couple canonique : rembobiner un prouveur pour obtenir deux défis sur un même
engagement extrait le témoin, ce qui est le sens de \og{} preuve de connaissance \fg{}, tandis que le
simulateur montre que le vérificateur n'a rien gagné --- la même forme d'argument qu'en CRY-17,
cette fois sur un groupe au lieu d'un graphe. Le lemme de bifurcation (forking lemma) de
Pointcheval et Stern (1996) a fait de l'étape de rembobinage un théorème général sur les
signatures. Le brevet de Schnorr a expiré en 2008, après quoi la construction s'est répandue
dans EdDSA et, en 2021, dans Bitcoin.

\begin{refs}
\item Contexte : Protocole d'authentification de Schnorr --- Wikipédia (\url{https://fr.wikipedia.org/wiki/Protocole_d%27authentification_de_Schnorr})
\item Contexte : Preuve de connaissance --- Wikipédia (en anglais) (\url{https://en.wikipedia.org/wiki/Proof_of_knowledge})
\item Article : C. P. Schnorr, \og{} Efficient signature generation by smart cards \fg{}, Journal of Cryptology 4 (1991)
\end{refs}

\bigskip


\section*{Recherche}

\begin{problem}[{Quelle est la sécurité des schémas post-quantiques à réseaux ? --- Recherche}]
\textbf{CRY-21.} \textbf{Problème.} En août 2024, le NIST a finalisé ses premières
normes post-quantiques : FIPS 203 (ML-KEM, issu de CRYSTALS-Kyber, encapsulation de clé),
FIPS 204 (ML-DSA, issu de CRYSTALS-Dilithium, signatures) --- tous deux fondés sur des réseaux
de modules et sur le problème Module-LWE --- et FIPS 205 (SLH-DSA, issu de SPHINCS+,
signatures fondées sur le hachage) ; en mars 2025 il a retenu HQC, fondé sur les codes, comme
mécanisme d'encapsulation de secours. La question mathématique sous-jacente est grande
ouverte. Comprenez l'état de l'art et repérez les points fragiles.
\begin{enumerate}
\item[(a)] Le problème (Module-)LWE
est-il réellement difficile --- classiquement ou quantiquement ? Aucun algorithme quantique
polynomial n'est connu, mais aucune démonstration n'en exclut l'existence ; la voie du
sous-groupe caché qui a cassé la factorisation ne s'applique pas, mais ce n'est qu'une preuve
d'absence, non une absence de preuve.
\item[(b)] Quelle est la finesse de la sécurité concrète ? Les meilleures attaques connues sont des
algorithmes de réduction de réseau et de criblage dont le temps est exponentiel en la
dimension ; les choix de paramètres extrapolent leur coût, et une véritable amélioration des
constantes de criblage rognerait des bits réels sur chaque clé déployée.
\item[(c)] La structure
algébrique supplémentaire des réseaux de modules et d'anneaux (sur des réseaux d'idéaux dans
des corps cyclotomiques) autorise-t-elle des attaques que le LWE non structuré ne permet
pas ? Des attaques partielles sur des problèmes voisins (par exemple sur certains cadres de
réseaux d'idéaux) sont connues ; rien n'atteint les schémas normalisés.
\end{enumerate}

\end{problem}

\noindent C'est la plus grande expérience en cours de la cryptographie : une
compétition publique d'une décennie (2016-2024) s'est achevée sur le début de la migration
mondiale de l'échange de clés vers les réseaux, sur la foi des réductions du pire cas de
Regev (CRY-16) et de vingt-cinq ans d'attaques infructueuses. Le statut honnête : aucune
attaque efficace, aucune démonstration, et des enjeux aussi élevés qu'il est possible.
L'histoire conseille l'humilité --- SIKE, finaliste fondé sur les isogénies, a été cassé par des
mathématiques classiques (Castryck--Decru) en 2022, quelques semaines après avoir été
retenu.

\begin{refs}
\item Contexte : Kyber --- Wikipédia (en anglais) (\url{https://en.wikipedia.org/wiki/Kyber})
\item Contexte : Cryptographie post-quantique --- Wikipédia (\url{https://fr.wikipedia.org/wiki/Cryptographie_post-quantique})
\item Norme : NIST FIPS 203, Module-Lattice-Based Key-Encapsulation Mechanism Standard (2024) (\url{https://csrc.nist.gov/pubs/fips/203/final})
\item Contexte : Annonce du NIST, août 2024 (\url{https://www.nist.gov/news-events/news/2024/08/nist-releases-first-3-finalized-post-quantum-encryption-standards})
\end{refs}

\bigskip

\begin{problem}[{Obfuscation indistinguable : de l'impossible à l'impraticable --- Recherche}]
\textbf{CRY-22.} \textbf{Problème.} Un obfuscateur indistinguable (iO)
compile les programmes de sorte que les obfuscations de deux circuits quelconques de même
taille calculant la même fonction soient calculatoirement indistinguables. iO est
\og{} crypto-complet \fg{} : avec lui (plus quelques hypothèses légères) on peut construire presque
tout objet cryptographique connu --- chiffrement fonctionnel, chiffrement déniable, et bien
d'autres. Statut, honnêtement : après le premier candidat (Garg et al., 2013) et les ruptures
successives de nombreux candidats à applications multilinéaires, Jain, Lin et Sahai ont
démontré (STOC 2021) que iO existe en supposant la difficulté sous-exponentielle de quatre
problèmes bien étudiés (couplages, LPN sur de grands corps, et générateurs pseudo-aléatoires
dans $\mathrm{NC}^0$). La construction n'est qu'un résultat de faisabilité --- astronomiquement
loin d'être praticable. Les points suivants sont ouverts ; comprenez l'architecture de
Jain--Lin--Sahai assez bien pour expliquer où chaque hypothèse intervient.
\begin{enumerate}
\item[(a)] Construire iO à partir d'une seule hypothèse standard,
idéalement LWE seul (l'obtention d'un iO post-quantique à partir d'hypothèses bien fondées
demeure non tranchée).
\item[(b)] Amener l'efficacité en vue d'une implémentation.
\item[(c)] Déterminer
si les hypothèses sous-exponentielles requises peuvent être affaiblies en hypothèses
polynomiales.
\end{enumerate}

\end{problem}

\noindent Pendant deux décennies, iO a oscillé entre \og{} manifestement impossible \fg{}
(l'obfuscation en boîte noire, elle, \textit{est} impossible --- Barak et al., 2001) et
\og{} manifestement trop beau pour être vrai \fg{}. Le théorème de 2021 a réglé l'existence sous des
hypothèses de type standard et il est considéré comme l'un des résultats les plus profonds de
l'histoire de la cryptographie ; la distance entre l'existence et le déploiement est
désormais la frontière.

\begin{refs}
\item Contexte : Obfuscation indistinguable --- Wikipédia (\url{https://fr.wikipedia.org/wiki/Obfuscation_indistinguable})
\item Article : A. Jain, H. Lin, A. Sahai, \og{} Indistinguishability Obfuscation from Well-Founded Assumptions \fg{} (STOC 2021), arXiv:2008.09317 (\url{https://arxiv.org/abs/2008.09317})
\end{refs}

\bigskip

\begin{problem}[{Quelle est la finesse des accélérations quantiques ? --- Recherche}]
\textbf{CRY-23.} \textbf{Problème.} L'algorithme de Grover fouille un espace non
structuré de $N$ objets en $O(\sqrt N)$ requêtes quantiques, et le théorème BBBV
(Bennett, Bernstein, Brassard, Vazirani, 1997) démontre que c'est optimal : l'accélération
quadratique est \textit{exactement} la limite pour la recherche en boîte noire, ce qui explique
que la cryptographie symétrique survive aux ordinateurs quantiques au prix d'un doublement des
longueurs de clés. Pour les problèmes structurés, le paysage est ouvert. Dressez la carte des
hypothèses cryptographiques qui tombent devant quelle classe d'attaque quantique, et de
l'endroit où les démonstrations s'arrêtent réellement.
\begin{enumerate}
\item[(a)] L'accélération exponentielle de Shor pour la factorisation se
mesure au meilleur algorithme classique \textit{connu} ; il n'existe aucune borne inférieure
classique surpolynomiale pour la factorisation, ni aucune démonstration que
$\mathrm{BQP} \ne \mathrm{BPP}$. Qu'impliquerait l'un ou l'autre énoncé ?
\item[(b)] Raz et
Tal (2019) ont démontré une séparation par oracle de BQP et de la hiérarchie polynomiale ---
dites soigneusement ce que les séparations par oracle établissent et n'établissent pas au
sujet du monde non relativisé.
\item[(c)] Statut honnête : pour les problèmes NP-complets, aucune
accélération quantique surpolynomiale n'est connue ni largement jugée possible (BBBV exclut
déjà la voie de la boîte noire) ; démontrer ou réfuter \og{} les ordinateurs quantiques ne
peuvent pas résoudre les problèmes NP-complets en temps polynomial \fg{} est aussi hors de portée
que P contre NP.
\end{enumerate}

\end{problem}

\noindent La cryptographie d'après Shor est un pari sur la structure fine de la
théorie de la complexité quantique. Les croyances de travail du domaine --- la factorisation est
facile quantiquement, LWE et AES-256 ne le sont pas --- reposent sur un théorème d'optimalité
démontré (BBBV), sur un algorithme (Shor) et sur des décennies d'échecs à étendre l'un ou
l'autre. Énoncer précisément ce qui est démontré, ce qui est conjecturé et ce qui n'est
qu'habitude est en soi un exercice de niveau recherche ; les plus profondes de ces questions
rejoignent les problèmes de complexité de la page
Problèmes ouverts (\url{open-problems.html}).

\begin{refs}
\item Contexte : Algorithme de Grover --- Wikipédia (\url{https://fr.wikipedia.org/wiki/Algorithme_de_Grover})
\item Contexte : BQP --- Wikipédia (\url{https://fr.wikipedia.org/wiki/BQP})
\item Article : C. Bennett, E. Bernstein, G. Brassard, U. Vazirani, \og{} Strengths and Weaknesses of Quantum Computing \fg{}, SIAM Journal on Computing (1997)
\item Article : R. Raz et A. Tal, \og{} Oracle Separation of BQP and PH \fg{}, STOC 2019
\end{refs}

\bigskip

\begin{problem}[{Les fonctions à sens unique existent-elles ? --- Recherche, short}]
\textbf{CRY-39.} \textbf{Problème.} Démontrez que si des fonctions à sens unique existent
alors $\mathrm{P} \ne \mathrm{NP}$, et expliquez précisément pourquoi l'implication
réciproque n'est pas connue. Énoncez ensuite, avec toutes les hypothèses exactes, le théorème
de Liu et Pass (2020) selon lequel des fonctions à sens unique existent si et seulement si la
complexité de Kolmogorov à temps borné $K^{t}$ est légèrement difficile en moyenne.
\end{problem}

\noindent \textbf{Statut : ouvert}, et c'est le fondement sur lequel repose tout le
reste de cette page : générateurs pseudo-aléatoires, chiffrement symétrique, signatures
numériques et mises en gage (CRY-35) existent tous si et seulement si les fonctions à sens
unique existent --- une seule phrase non démontrée porte donc la cryptographie à clé secrète
tout entière. Démontrer qu'elles existent réglerait $\mathrm{P} \ne \mathrm{NP}$ en
corollaire, ce qui explique que personne n'attende de démonstration prochaine. Liu et Pass ont
changé la forme de la question en trouvant, après quarante ans de recherche, un problème
\textit{naturel} en moyenne équivalent à leur existence --- calculer le plus court programme qui
produit une chaîne donnée dans une limite de temps.

\begin{refs}
\item Contexte : Fonction à sens unique --- Wikipédia (\url{https://fr.wikipedia.org/wiki/Fonction_%C3%A0_sens_unique})
\item Article : Y. Liu \& R. Pass, \og{} On One-way Functions and Kolmogorov Complexity \fg{}, FOCS 2020
\end{refs}

\bigskip

\begin{problem}[{L'écart d'approximation des réseaux --- Recherche, short}]
\textbf{CRY-40.} \textbf{Problème.} L'algorithme LLL (1982) produit, en temps polynomial,
un vecteur non nul du réseau qui n'est pas plus long que $2^{(n-1)/2}$ fois le plus court,
tandis que le problème du plus court vecteur est NP-difficile sous réductions aléatoires pour
tout facteur d'approximation constant, et qu'une approximation à $\sqrt{n}$ près appartient
à $\mathrm{NP} \cap \mathrm{coNP}$ et est donc peu susceptible d'être NP-difficile.
Déterminez où se trouve la vérité : trouvez le plus grand facteur d'approximation pour lequel
le problème est difficile, ou le plus petit pour lequel il est facile.
\end{problem}

\noindent \textbf{Statut : ouvert}, avec un écart exponentiellement large --- le meilleur
facteur polynomial connu en temps polynomial reste seulement $2^{O(n \log\log n / \log n)}$,
à peine en deçà de celui de LLL, et aucun facteur d'approximation polynomial n'est atteignable
en temps polynomial par une méthode connue. La cryptographie post-quantique (CRY-16, CRY-21)
est déployée à l'intérieur de cet écart : sa sécurité suppose que le problème reste difficile
pour des facteurs polynomiaux en $n$, région où n'existent ni démonstration de difficulté ni
algorithme. LLL fut lui-même un produit dérivé --- Lenstra, Lenstra et Lovász l'ont inventé pour
factoriser les polynômes à coefficients rationnels, et ce n'est que plus tard qu'il est devenu
l'instrument standard de la cryptanalyse des réseaux.

\begin{refs}
\item Contexte : Problème de réseau --- Wikipédia (\url{https://fr.wikipedia.org/wiki/Probl%C3%A8me_de_r%C3%A9seau})
\item Contexte : Algorithme de réduction de base de réseau de Lenstra--Lenstra--Lovász --- Wikipédia (\url{https://fr.wikipedia.org/wiki/Algorithme_LLL})
\item Article : A. K. Lenstra, H. W. Lenstra Jr. \& L. Lovász, \og{} Factoring polynomials with rational coefficients \fg{}, Mathematische Annalen 261 (1982)
\item Article : D. Aharonov \& O. Regev, \og{} Lattice problems in NP $\cap$ coNP \fg{}, Journal of the ACM 52 (2005)
\end{refs}

\bigskip

\begin{problem}[{La transformation de Fiat--Shamir est-elle solide hors du modèle de l'oracle aléatoire ? --- Recherche}]
\textbf{CRY-49.} \textbf{Problème.} La transformation de Fiat--Shamir rend non interactive une preuve
interactive à pièces publiques en remplaçant le défi aléatoire du vérificateur par un haché de
tout ce qui a été envoyé jusque-là : $c = H(x, a)$. Si $H$ est modélisée par un \textit{oracle
aléatoire} --- une fonction uniformément aléatoire que toute partie, honnête ou non, ne peut
évaluer qu'en l'interrogeant --- la transformation se démontre solide, et pratiquement tout
schéma de signature déployé de type Schnorr (CRY-48) repose sur une telle démonstration. La
question est de savoir ce qui subsiste quand $H$ est une fonction publiée réelle.
\begin{enumerate}
\item[(a)] Comprenez les résultats négatifs. Canetti, Goldreich et Halevi (1998) ont exhibé des
schémas démontrablement sûrs dans le modèle de l'oracle aléatoire qui sont non sûrs pour
\textit{toute} instanciation concrète de l'oracle ; Goldwasser et Kalai (2003) ont fait de même
spécifiquement pour Fiat--Shamir, en construisant un argument interactif dont l'effondrement à
la Fiat--Shamir est non sûr quelle que soit la fonction de hachage. Énoncez précisément les
deux résultats, puis dites exactement ce qu'ils excluent et n'excluent pas au sujet des
signatures de Schnorr.
\item[(b)] Comprenez les résultats positifs. Une famille de fonctions de hachage est \textit{intraitable
en corrélation} pour une relation si aucun adversaire efficace ne peut trouver un $x$ tel
que $(x, H(x))$ y appartienne ; pour des relations convenables, cette propriété rend
Fiat--Shamir solide sans aucune idéalisation, et de telles familles ont été construites à
partir de l'apprentissage avec erreurs (CRY-16) par Canetti, Chen, Holmgren, Lombardi,
Rothblum, Rothblum et Wichs en 2019, donnant des preuves non interactives à divulgation nulle
pour NP sous une hypothèse de réseau standard. Déterminez quels protocoles cette machinerie
couvre et lesquels elle laisse intacts.
\item[(c)] Le problème ouvert : démontrez ou réfutez que le schéma de signature Fiat--Shamir--Schnorr,
avec une fonction de hachage concrète, est infalsifiable sous une hypothèse standard
réfutable. Déterminez également si la perte de sécurité inhérente au rembobinage --- le lemme de
bifurcation abandonne un facteur égal au nombre de requêtes de hachage --- est nécessaire ou
n'est qu'un artefact de la technique de démonstration.
\end{enumerate}
\textit{Statut : (a) et (b) sont des théorèmes ; (c) est ouvert. Aucune démonstration dans le
modèle standard n'est connue pour le schéma de signature de Schnorr tel qu'il est déployé, et
des résultats de méta-réduction (Paillier et Vergnaud, 2005) montrent que des familles
entières de stratégies de démonstration naturelles ne peuvent pas aboutir.}
\end{problem}

\noindent Bellare et Rogaway ont proposé le modèle de l'oracle aléatoire en 1993
comme compromis délibéré : démontrer la sécurité contre un haché idéalisé, puis l'instancier
et espérer. Trente ans après, le compromis tient toujours et reste toujours injustifié --- les
schémas qui signent le trafic du monde ont des démonstrations dans un modèle dont on sait, par
la question (a), qu'il n'est pas solide en général, et pourtant aucun contre-exemple n'a jamais
touché un schéma que quiconque utilise. La version quantique de la question a été réglée par
l'affirmative en 2019, quand Don, Fehr, Majenz et Schaffner et, indépendamment, Liu et Zhandry
ont montré que Fiat--Shamir survit à un adversaire disposant d'un accès quantique à l'oracle.
La distance entre \og{} démontrable dans un modèle idéalisé \fg{} et \og{} utilisé par tous les
navigateurs \fg{} est la plus grosse dette impayée de la cryptographie appliquée.

\begin{refs}
\item Contexte : Heuristique de Fiat--Shamir --- Wikipédia (\url{https://fr.wikipedia.org/wiki/Heuristique_de_Fiat-Shamir})
\item Contexte : Modèle de l'oracle aléatoire --- Wikipédia (\url{https://fr.wikipedia.org/wiki/Mod%C3%A8le_de_l%27oracle_al%C3%A9atoire})
\item Article : S. Goldwasser et Y. Tauman Kalai, \og{} On the (In)security of the Fiat-Shamir Paradigm \fg{}, FOCS 2003
\item Article : R. Canetti, Y. Chen, J. Holmgren, A. Lombardi, G. N. Rothblum, R. D. Rothblum et D. Wichs, \og{} Fiat-Shamir: From Practice to Theory \fg{}, STOC 2019
\end{refs}

\bigskip

\begin{problem}[{Le transfert inconscient peut-il se construire à partir de fonctions à sens unique ? --- Recherche, short}]
\textbf{CRY-50.} \textbf{Problème.} Dans le \textit{transfert inconscient} 1 parmi 2, un émetteur
détient deux chaînes $m_0, m_1$ et un récepteur détient un bit $b$ ; à la fin du protocole
le récepteur connaît $m_b$ et rien de $m_{1-b}$, tandis que l'émetteur n'a rien appris
sur $b$. Démontrez que le transfert inconscient implique la mise en accord de clé, puis
tranchez le sens réciproque de la hiérarchie : décidez si le transfert inconscient peut être
construit à partir des seules fonctions à sens unique (CRY-39).
\end{problem}

\noindent \textbf{Statut : ouvert, et jugé impossible par toute voie en boîte noire.}
Michael Rabin a proposé le transfert inconscient en 1981, et Even, Goldreich et Lempel en ont
donné la forme 1 parmi 2 en 1985 ; Joe Kilian a démontré en 1988 qu'il est \textit{complet} ---
disposant du transfert inconscient, des parties mutuellement méfiantes peuvent calculer de
façon sûre n'importe quelle fonction --- ce qui en fait le pivot de toute la théorie du calcul
multipartite. La réduction demandée ci-dessus le place au-dessus de la mise en accord de clé
dans la carte des mondes possibles d'Impagliazzo, et Impagliazzo et Rudich ont montré en 1989
que la mise en accord de clé n'admet aucune construction en boîte noire à partir de fonctions
à sens unique : une construction classique exigerait donc des techniques que personne ne
possède. Le développement récent le plus frappant contourne la question tout entière :
Bartusek, Coladangelo, Khurana et Ma et, indépendamment, Grilo, Lin, Song et Vaikuntanathan
ont montré en 2021 que si les parties peuvent échanger des états quantiques, alors les
fonctions à sens unique \textit{suffisent} --- la séparation classique est donc un énoncé sur la
communication classique, non sur la cryptographie.

\begin{refs}
\item Contexte : Transfert inconscient --- Wikipédia (\url{https://fr.wikipedia.org/wiki/Transfert_inconscient})
\item Contexte : Calcul multipartite sécurisé --- Wikipédia (\url{https://fr.wikipedia.org/wiki/Calcul_multipartite_s%C3%A9curis%C3%A9})
\item Article : J. Kilian, \og{} Founding Cryptography on Oblivious Transfer \fg{}, STOC 1988
\item Article : A. B. Grilo, H. Lin, F. Song et V. Vaikuntanathan, \og{} Oblivious Transfer is in MiniQCrypt \fg{} (2020), arXiv:2011.14980 (\url{https://arxiv.org/abs/2011.14980})
\end{refs}

\bigskip


\section*{Pour aller plus loin}


\vfill
\noindent\emph{Hors de la Caverne} --- des probl\`emes, pas de r\'eponses.
\end{document}
